\input{preamble}

% OK, start here.
%
\begin{document}

\title{La formule des traces}


\maketitle

\phantomsection
\label{section-phantom}

\tableofcontents


\section{Introduction}
\label{section-introduction}

\noindent
Voici les notes de la seconde partie d'un cours de cohomologie \'etale
donné par Johan de Jong à l'Université Columbia à l'automne 2009. Les
auteurs des notes initiales étaient Thibaut Pugin, Zachary Maddock et Min Lee.
Au fil du temps, nous ajouterons des références aux prérequis dans le reste du
projet Champs et donnerons des démonstrations rigoureuses de tous les énoncés.









\section{La formule des traces}
\label{section-trace-formula}

\noindent
Un cours classique de cohomologie \'etale énoncerait et démontrerait normalement
les théorèmes de changement de base propre et lisse, la pureté et la dualité de
Poincar\'e. Tout cela se trouve dans \cite[Arcata]{SGA4.5}. Nous allons plutôt
étudier la formule des traces pour le Frobenius, en suivant l'exposé de Deligne
dans \cite[Rapport]{SGA4.5}. Nous nous limiterons à la dimension 1, mais le
changement de base propre suffit alors à traiter le cas général. Puisque tous
les groupes de cohomologie considérés seront \'etales, nous omettons l'indice
$_\etale$. Décrivons maintenant
la formule que nous cherchons. Soient $X$ un schéma de type fini de dimension 1
sur un corps fini $k$, $\ell$ un nombre premier et $\mathcal{F}$ un faisceau
constructible et plat de $\mathbf{Z}/\ell^n\mathbf{Z}$-modules. Alors
\begin{equation}
\label{equation-trace-formula-initial}
\sum\nolimits_{x \in X(k)}
\text{Tr}(\text{Frob} | \mathcal{F}_{\bar x}) =
\sum\nolimits_{i = 0}^2
(-1)^i \text{Tr}(\pi_X^* | H^i_c(X \otimes_k \bar k, \mathcal{F}))
\end{equation}
comme éléments de $\mathbf{Z}/\ell^n\mathbf{Z}$. Comme nous le verrons, cette
formulation est légèrement incorrecte telle quelle. Décrivons néanmoins les
symboles qui y figurent.




\section{Les Frobenius}
\label{section-frobenii}

\noindent
Dans cette section, nous allons démontrer un théorème ``déroutant''.
Un analogue topologique de ce théorème déroutant est le suivant.

\begin{exercise}
\label{exercise-baffling}
Soient $X$ un espace topologique et $g : X \to X$ une application continue telle
que $g^{-1}(U) = U$ pour tout ouvert $U$ de $X$. Alors $g$ induit l'identité sur
la cohomologie de $X$ (à coefficients quelconques).
\end{exercise}

\noindent
Passons maintenant à l'énoncé pour le site \'etale.

\begin{lemma}
\label{lemma-baffling}
Soient $X$ un schéma et $g : X \to X$ un morphisme. Supposons que, pour tout
morphisme \'etale $\varphi : U \to X$, il existe un isomorphisme
$$
\xymatrix{
U \ar[rd]_\varphi \ar[rr]^-\sim & & {U
\times_{\varphi, X, g} X} \ar[ld]^{\text{pr}_2} \\
& X
}
$$
fonctoriel en $U$. Alors $g$ induit l'identité en cohomologie (pour tout faisceau).
\end{lemma}

\begin{proof}
La démonstration est formelle et ne présente aucune difficulté.
\end{proof}

\noindent
Voir Variétés, section \ref{varieties-section-frobenius},
pour une discussion des différentes variantes du morphisme de Frobenius.

\begin{theorem}[Le théorème déroutant]
\label{theorem-baffling}
Soit $X$ un schéma de caractéristique $p > 0$. Alors le Frobenius absolu
induit (par image réciproque) l'identité en cohomologie, c'est-à-dire que, pour
tout entier $j\geq 0$,
$$
F_X^* : H^j (X, \underline{\mathbf{Z}/n\mathbf{Z}}) \longrightarrow H^j (X,
\underline{\mathbf{Z}/n\mathbf{Z}})
$$
est l'identité.
\end{theorem}

\noindent
Ce théorème est purement formel. Il est toutefois utile de revoir comment
calculer l'image réciproque d'une classe de cohomologie. Disons simplement que,
lorsque la cohomologie coïncide avec celle de {\v C}ech, il suffit de tirer en
arrière les cocycles de {\v C}ech (au moyen des produits fibrés sur un site). Le
cas général est assez technique ; voir
Hyperrecouvrements, théorème \ref{hypercovering-theorem-cohomology-hypercoverings}.
Pour démontrer le théorème, il nous suffit
de vérifier que l'hypothèse du lemme \ref{lemma-baffling}
est satisfaite par le Frobenius.

\begin{proof}[Démonstration du théorème \ref{theorem-baffling}]
Nous devons vérifier l'existence d'un isomorphisme fonctoriel comme ci-dessus.
Pour un morphisme \'etale $\varphi : U \to X$, considérons le diagramme
$$
\xymatrix{
U \ar@{-->}[rd] \ar@/^1pc/[rrd]^{F_U}
\ar@/_1pc/[rdd]_\varphi \\
& {U \times_{\varphi, X, F_X} X} \ar[r]_-{\text{pr}_1}
\ar[d]^{\text{pr}_2} & U \ar[d]^\varphi \\
& X \ar[r]^{F_X} & X.
}
$$
La flèche en pointillés est un morphisme \'etale et un homéomorphisme
universel ; c'est donc un isomorphisme. Voir
Morphismes étales, lemme \ref{etale-lemma-relative-frobenius-etale}.
\end{proof}

%10.22.09

\begin{definition}
\label{definition-geometric-frobenius}
Soit $k$ un corps fini à $q = p^f$ éléments. Soit $X$ un schéma
sur $k$. Le {\it Frobenius géométrique} de $X$ est le morphisme
$\pi_X : X \to X$ au-dessus de $\Spec(k)$ égal à $F_X^f$.
\end{definition}

\noindent
Puisque $\pi_X$ est un morphisme sur $k$, nous pouvons le changer de base par
n'importe quel schéma sur $k$. En particulier, le changement de base à la
clôture algébrique $\bar k$ donne un morphisme $\pi_X : X_{\bar k} \to
X_{\bar k}$. Employer aussi $\pi_X$ pour ce changement de base ne devrait pas
prêter à confusion, car $X_{\bar k}$ n'a pas de Frobenius géométrique propre.

\begin{lemma}
\label{lemma-sheaf-over-finite-field-has-frobenius-descent}
Soit $\mathcal{F}$ un faisceau sur $X_\etale$.
Il existe alors des isomorphismes canoniques
$\pi_X^{-1} \mathcal{F} \cong \mathcal{F}$ et
$\mathcal{F} \cong {\pi_X}_*\mathcal{F}$.
\end{lemma}

\noindent
Cela est faux pour le site fppf.

\begin{proof}
Soit $\varphi : U \to X$ un morphisme \'etale. Rappelons que
${\pi_X}_* \mathcal{F} (U) = \mathcal{F} (U \times_{\varphi, X, \pi_X} X)$.
Puisque $\pi_X = F_X^f$, la démonstration du
théorème \ref{theorem-baffling} fournit un isomorphisme fonctoriel
$$
\xymatrix{
U \ar[rd]_{\varphi} \ar[rr]_-{\gamma_U}
& & U \times_{\varphi, X, \pi_X} X \ar[ld]^{\text{pr}_2} \\
& X
}
$$
où $\gamma_U = (\varphi, F_U^f)$. Nous définissons maintenant un
isomorphisme
$$
\mathcal{F} (U) \longrightarrow {\pi_X}_* \mathcal{F} (U) =
\mathcal{F} (U \times_{\varphi, X, \pi_X} X)
$$
comme l'application de restriction de $\mathcal{F}$ le long de $\gamma_U^{-1}$.
L'autre isomorphisme se construit de manière analogue.
\end{proof}

\begin{remark}
\label{remark-may-be-confusing}
Il se peut, ou non, que $F^f_U$ soit égal à $\pi_U$.
\end{remark}

\noindent
Poursuivons l'étude de la cohomologie des faisceaux sur notre schéma $X$ défini
sur le corps fini $k$ à $q = p^f$ éléments.
Fixons une clôture algébrique $\bar k$ de $k$ et notons $G_k =
\text{Gal}(\bar k/k)$ le groupe de Galois absolu de $k$.
Soit $\mathcal{F}$ un faisceau abélien sur $X_\etale$.
Nous allons munir d'une structure de $G_k$-module à gauche le
groupe de cohomologie $H^j (X_{\bar k}, \mathcal{F}|_{X_{\bar k}})$
de la manière suivante : si $\sigma \in G_k$, le diagramme
$$
\xymatrix{
X_{\bar k} \ar[rd] \ar[rr]^{\Spec(\sigma) \times \text{id}_X} & &
X_{\bar k} \ar[ld] \\
& X
}
$$
est commutatif. Nous pouvons donc poser, pour $\xi \in H^j (X_{\bar k},
\mathcal{F}|_{X_{\bar k}})$,
$$
\sigma \cdot \xi := (\Spec(\sigma) \times \text{id}_X)^*\xi \in
H^j(X_{\bar k}, (\Spec(\sigma) \times \text{id}_X)^{-1}
\mathcal{F}|_{X_{\bar k}})
= H^j (X_{\bar k}, \mathcal{F}|_{X_{\bar k}}),
$$
où la dernière égalité résulte de la commutativité du diagramme précédent.
Cela munit ce dernier groupe d'une structure de $G_k$-module.

\begin{lemma}
\label{lemma-two-actions-agree}
Dans la situation ci-dessus, notons $\alpha : X \to \Spec(k)$ le morphisme
structural. Considérons la fibre $(R^j\alpha_*\mathcal{F})_{\Spec(\bar k)}$
munie de son action galoisienne naturelle comme dans
Cohomologie étale, section \ref{etale-cohomology-section-galois-action-stalks}.
Alors l'identification
$$
(R^j\alpha_*\mathcal{F})_{\Spec(\bar k)} \cong H^j (X_{\bar k},
\mathcal{F}|_{X_{\bar k}})
$$
du
théorème \ref{etale-cohomology-theorem-higher-direct-images} de Cohomologie étale
est un isomorphisme de $G_k$-modules.
\end{lemma}

\noindent
Un résultat analogue vaut lorsque l'on compare
$(R^j\alpha_!\mathcal{F})_{\Spec(\bar k)}$ à
$H^j_c (X_{\bar k}, \mathcal{F}|_{X_{\bar k}})$.

\begin{proof}
Omise.
\end{proof}

\begin{definition}
\label{definition-arithmetic-frobenius}
Le {\it Frobenius arithmétique} est l'application
$\text{frob}_k : \bar k \to \bar k$, $x \mapsto x^q$ appartenant à $G_k$.
\end{definition}

\begin{theorem}
\label{theorem-geometric-arithmetic-inverse}
Soit $\mathcal{F}$ un faisceau abélien sur $X_\etale$. Alors, pour tout
$j\geq 0$, $\text{frob}_k$ agit sur le groupe de cohomologie $H^j(X_{\bar k},
\mathcal{F}|_{X_{\bar k}})$ comme l'inverse de l'application $\pi_X^*$.
\end{theorem}

\noindent
L'application $\pi_X^*$ est définie par la composée
$$
H^j(X_{\bar k}, \mathcal{F}|_{X_{\bar k}}) \xrightarrow{{\pi_X}_{\bar k}^*}
H^j(X_{\bar k}, (\pi_X^{-1} \mathcal{F})|_{X_{\bar k}}) \cong
H^j(X_{\bar k}, \mathcal{F}|_{X_{\bar k}}).
$$
où le dernier isomorphisme provient de l'isomorphisme canonique
$\pi_X^{-1} \mathcal{F} \cong \mathcal{F}$ du
lemme \ref{lemma-sheaf-over-finite-field-has-frobenius-descent}.

\begin{proof}
La composée $X_{\bar k} \xrightarrow{\Spec(\text{frob}_k)} X_{\bar k}
\xrightarrow{\pi_X} X_{\bar k}$ est égale à $F_{X_{\bar k}}^f$ ; le
résultat découle donc du théorème déroutant convenablement généralisé à des
coefficients non triviaux. Notons que la composée précédente commute au sens où
$F_{X_{\bar k}}^f = \pi_X \circ \Spec(\text{frob}_k) =
\Spec(\text{frob}_k) \circ \pi_X$.
\end{proof}

\begin{definition}
\label{definition-geometric-frobenius-on-stalk}
Si $x \in X(k)$ est un point rationnel et $\bar x : \Spec(\bar k) \to X$
est le point géométrique au-dessus de $x$, nous notons $\pi_x : \mathcal{F}_{\bar x} \to
\mathcal{F}_{\bar x}$ l'action de $\text{frob}_k^{-1}$ et l'appelons le
{\it Frobenius géométrique}\footnote{Cette notation n'est pas standard.
Cet opérateur est noté $F_x$ dans \cite{SGA4.5}. Nous changerons probablement
cette notation à l'avenir.}
\end{definition}

\noindent
Nous pouvons maintenant donner un énoncé plus précis (quoique faux) de la formule des
traces (\ref{equation-trace-formula-initial}). Soit $X$ un schéma de
type fini de dimension 1
sur un corps fini $k$, $\ell$ un nombre premier et $\mathcal{F}$ un faisceau
constructible et plat de $\mathbf{Z}/\ell^n\mathbf{Z}$-modules. Alors
\begin{equation}
\label{equation-trace-formula-second}
\sum\nolimits_{x \in X(k)}
\text{Tr}(\pi_x | \mathcal{F}_{\bar x})
=
\sum\nolimits_{i = 0}^2
(-1)^i \text{Tr}(\pi_X^* | H^i_c(X_{\bar k}, \mathcal{F}))
\end{equation}
comme éléments de $\mathbf{Z}/\ell^n\mathbf{Z}$. Cette égalité est incorrecte
parce que la trace du membre de droite n'a pas de sens pour le type de
faisceaux considéré. Avant d'aborder ce problème, cherchons à motiver
l'apparition du Frobenius géométrique (outre le fait qu'il s'agit d'un
morphisme naturel !).

\medskip\noindent
Considérons le cas où $X = \mathbf{P}^1_k$ et $\mathcal{F} =
\underline{\mathbf{Z}/\ell\mathbf{Z}}$. En tout point, le module galoisien
$\mathcal{F}_{\bar x}$ est trivial ; ainsi, pour tout morphisme $\varphi$ agissant sur
$\mathcal{F}_{\bar x}$, le membre de gauche vaut
$$
\sum\nolimits_{x \in X(k)} \text{Tr}(\varphi | \mathcal{F}_{\bar x}) =
\#\mathbf{P}^1_k(k) = q+1.
$$
Or $\mathbf{P}^1_k$ est propre, donc la cohomologie à support propre coïncide avec la
cohomologie usuelle ; ainsi, pour un morphisme $\pi : \mathbf{P}^1_k \to
\mathbf{P}^1_k$, le membre de droite vaut
$$
\text{Tr}(\pi^* | H^0 (\mathbf{P}^1_{\bar k},
\underline{\mathbf{Z}/\ell\mathbf{Z}})) + \text{Tr}(\pi^* | H^2
(\mathbf{P}^1_{\bar k}, \underline{\mathbf{Z}/\ell\mathbf{Z}})).
$$
Le module galoisien $H^0 (\mathbf{P}^1_{\bar k},
\underline{\mathbf{Z}/\ell\mathbf{Z}}) = \mathbf{Z}/\ell\mathbf{Z}$ est trivial,
car l'image réciproque de l'identité est l'identité. La première trace vaut donc 1,
quel que soit $\pi$. Pour la seconde trace, nous devons calculer l'image réciproque
$\pi^* : H^2(\mathbf{P}^1_{\bar k}, \underline{\mathbf{Z}/\ell\mathbf{Z}}))$
pour une application $\pi : \mathbf{P}^1_{\bar k} \to \mathbf{P}^1_{\bar k}$. C'est un
bon exercice, et la réponse est la multiplication par le degré de $\pi$
(pour une démonstration, voir
Cohomologie étale, lemme \ref{etale-cohomology-lemma-pullback-on-h2-curve}).
Autrement dit, le calcul est le même que dans le cadre familier de la cohomologie complexe.
En particulier, si $\pi$ est le Frobenius géométrique, nous obtenons
$$
\text{Tr}(\pi_X^* | H^2 (\mathbf{P}^1_{\bar k},
\underline{\mathbf{Z}/\ell\mathbf{Z}})) = q
$$
et si $\pi$ est le Frobenius arithmétique, nous obtenons
$$
\text{Tr}(\text{frob}_k^* | H^2 (\mathbf{P}^1_{\bar k},
\underline{\mathbf{Z}/\ell\mathbf{Z}})) = q^{-1}.
$$
Cette dernière possibilité est manifestement incorrecte.

\begin{remark}
\label{remark-compute-degree-lifting}
Le calcul des degrés peut s'effectuer en relevant (en un sens évident)
en caractéristique 0 et en considérant la situation à coefficients complexes.
Cette méthode ne fonctionne presque jamais, puisque le relèvement est en général impossible pour les
schémas autres que l'espace projectif.
\end{remark}

\noindent
Reste à comprendre pourquoi nous devons considérer la cohomologie à support
propre. En fait, au vu de la dualité de Poincar\'e, elle n'est pas absolument
nécessaire pour les variétés lisses, mais cela impose d'introduire certaines puissances
de $q$. Par exemple, considérons le cas où
$X = \mathbf{A}^1_k$ et
$\mathcal{F} = \underline{\mathbf{Z}/\ell\mathbf{Z}}$.
L'action sur les fibres est encore triviale ; nous n'avons donc qu'à examiner l'action
sur la cohomologie. Mais alors $\pi_X^*$ agit comme l'identité sur
$H^0(\mathbf{A}^1_{\bar k}, \underline{\mathbf{Z}/\ell\mathbf{Z}})$
et comme la multiplication par $q$ sur
$H^2_c(\mathbf{A}^1_{\bar k}, \underline{\mathbf{Z}/\ell\mathbf{Z}})$.




\section{Traces}
\label{section-traces}

\noindent
Nous expliquons maintenant comment prendre la trace d'un endomorphisme d'un module sur un
anneau non commutatif. Fixons un anneau fini $\Lambda$ dont le cardinal est premier à $p$.
En général, $\Lambda$ est l'anneau de groupe $(\mathbf{Z}/\ell^n\mathbf{Z})[G]$ d'un
groupe fini $G$. Par convention, tous les $\Lambda$-modules considérés
seront des $\Lambda$-modules à gauche.

\medskip\noindent
Introduisons la notation suivante :
nous notons $\Lambda^\natural$ le quotient de $\Lambda$ par le sous-groupe additif
engendré par les commutateurs (c'est-à-dire les éléments de la forme
$ab-ba$, $a, b \in \Lambda$). Notons que $\Lambda^\natural$ n'est pas un anneau.

\medskip\noindent
Par exemple, le module $(\mathbf{Z}/\ell^n\mathbf{Z})[G]^\natural$ est le
dual de l'espace des fonctions centrales, de sorte que
$$
(\mathbf{Z}/\ell^n\mathbf{Z})[G]^\natural
=
\bigoplus\nolimits_{\text{classes de conjugaison dans }G}
\mathbf{Z}/\ell^n\mathbf{Z}.
$$
Pour un $\Lambda$-module libre, on a $\text{End}_\Lambda(\Lambda^{\oplus m}) =
\text{Mat}_m(\Lambda)$. Puisque les modules sont à gauche,
la représentation des endomorphismes par des matrices est une action à droite : si $a \in
\text{End}(\Lambda^{\oplus m})$ a pour matrice $A$ et si $v \in \Lambda$, alors $a(v)
= v A$.

\begin{definition}
\label{definition-trace}
La {\it trace} de l'endomorphisme $a$ est la somme des termes diagonaux d'une
matrice qui le représente. On obtient ainsi une application additive $\text{Tr} :
\text{End}_\Lambda(\Lambda^{\oplus m}) \to \Lambda^\natural$.
\end{definition}

\begin{exercise}
\label{exercise-trace-is-trace}
Étant données des applications
$$
\Lambda^{\oplus m} \xrightarrow{a} \Lambda^{\oplus n}
\quad\text{et}\quad
\Lambda^{\oplus n} \xrightarrow{b} \Lambda^{\oplus m}
$$
montrer que $\text{Tr}(ab) = \text{Tr}(ba)$.
\end{exercise}

\noindent
Étendons la définition de la trace à un $\Lambda$-module projectif de type fini
$P$ et à un endomorphisme $\varphi$ de $P$. Écrivons $P$ comme facteur direct
d'un $\Lambda$-module libre, c'est-à-dire considérons des applications $P \xrightarrow{a}
\Lambda^{\oplus n} \xrightarrow{b} P$ telles que
\begin{enumerate}
\item
$\Lambda^{\oplus n} = \Im(a) \oplus \Ker(b)$ ; et
\item
$b\circ a = \text{id}_P$.
\end{enumerate}
Nous posons alors $\text{Tr}(\varphi) = \text{Tr}(a\varphi b)$. Il est facile de vérifier
que cette définition ne dépend pas des choix, au moyen de l'exercice précédent.








\section{Pourquoi les catégories dérivées ?}
\label{section-derived-categories-why}

\noindent
Avec cette définition de la trace, examinons un autre problème que pose la
formule telle qu'elle est énoncée. Soit $C$ une courbe projective lisse sur $k$. Il existe
une correspondance entre les faisceaux finis localement constants $\mathcal{F}$ sur
$C_\etale$ dont les fibres sont isomorphes à
${(\mathbf{Z}/\ell^n\mathbf{Z})}^{\oplus m}$ d'une part, et les représentations continues
$\rho : \pi_1 (C, \bar c) \to
\text{GL}_m(\mathbf{Z}/\ell^n\mathbf{Z}))$ (pour un choix fixé de $\bar c$)
d'autre part. Nous notons $\mathcal{F}_\rho$ le faisceau correspondant à
$\rho$. Alors $H^2 (C_{\bar k}, \mathcal{F}_\rho)$ est le groupe des coinvariants
pour l'action de $\rho(\pi_1 (C, \bar c))$ sur
${(\mathbf{Z}/\ell^n\mathbf{Z})}^{\oplus m}$, et l'on a la suite exacte courte
suivante
$$
0 \longrightarrow \pi_1 (C_{\bar k}, \bar c) \longrightarrow \pi_1 (C, \bar c)
\longrightarrow G_k \longrightarrow 0.
$$
Par exemple, faisons agir $\mathbf{Z} = \mathbf{Z} \sigma$ sur
$\mathbf{Z}/\ell^2\mathbf{Z}$ par $\sigma(x) = (1+\ell) x$. Les coinvariants
sont $(\mathbf{Z}/\ell^2\mathbf{Z})_{\sigma} = \mathbf{Z}/\ell\mathbf{Z}$, qui
n'est pas un $\mathbf{Z}/\ell^2\mathbf{Z}$-module plat. Nous ne pouvons donc pas prendre la
trace d'une action sur $H^2(C_{\bar k}, \mathcal{F}_\rho)$, du moins pas au
sens de la section précédente.

\medskip\noindent
En fait, notre objectif est d'établir une formule des traces à coefficients $\ell$-adiques.
Mais $\mathbf{Q}_\ell = \mathbf{Z}_\ell[1/\ell]$ et $\mathbf{Z}_\ell =
\lim \mathbf{Z}/\ell^n\mathbf{Z}$, et même pour un faisceau plat de
$\mathbf{Z}/\ell^n\mathbf{Z}$-modules, les groupes de cohomologie individuels peuvent ne pas
être plats, de sorte que nous ne pouvons calculer les traces. Une solution possible consiste à considérer le complexe
dérivé total $R\Gamma(C_{\bar k}, \mathcal{F}_\rho)$ dans la catégorie dérivée
$D(\mathbf{Z}/\ell^n\mathbf{Z})$ et à montrer qu'il s'agit d'un objet parfait, c'est-à-dire
qu'il est quasi-isomorphe à un complexe fini de modules libres de type fini.
Pour de tels complexes, on peut définir la trace, mais cela exigera un exposé
des catégories dérivées.






\section{Catégories dérivées}
\label{section-derived-categories}

\noindent
Pour fixer les notations, soit $\mathcal{A}$ une catégorie abélienne. Soit
$\text{Comp}(\mathcal{A})$ la catégorie abélienne des complexes de
$\mathcal{A}$. Soit $K(\mathcal{A})$ la catégorie des complexes à
homotopie près, dont les objets sont les complexes de $\mathcal{A}$ et les morphismes
des classes
d'homotopie de morphismes de complexes. Ce n'est pas une catégorie abélienne.
Grossièrement, $D(A)$ est définie comme la catégorie obtenue en inversant
tous les quasi-isomorphismes de $\text{Comp}(\mathcal{A})$ ou, de manière équivalente, de
$K(\mathcal{A})$. On peut en outre définir $\text{Comp}^+(\mathcal{A}),
K^+(\mathcal{A}), D^+(\mathcal{A})$ de manière analogue en ne considérant que des complexes bornés inférieurement.
De même, on peut définir $\text{Comp}^-(\mathcal{A}),
K^-(\mathcal{A}), D^-(\mathcal{A})$ à l'aide des complexes bornés supérieurement, et l'on peut
définir $\text{Comp}^b(\mathcal{A}), K^b(\mathcal{A}), D^b(\mathcal{A})$ à l'aide des
complexes bornés.

\begin{remark}
\label{remarks-derived-categories}
Quelques remarques sur les catégories dérivées.
\begin{enumerate}
\item
Des problèmes ensemblistes se posent lorsque $\mathcal{A}$ est assez
arbitraire ; nous les ignorerons allègrement.
\item
Les catégories $K(A)$ et $D(A)$ sont munies d'une structure de
catégorie triangulée.
\item
Les catégories $\text{Comp}(\mathcal{A})$ et $K(\mathcal{A})$ peuvent aussi être
définies lorsque $\mathcal{A}$ est une catégorie additive.
\end{enumerate}
\end{remark}

\noindent
Le foncteur d'homologie $H^i : \text{Comp}(\mathcal{A}) \to \mathcal{A}$ défini par
$K^\bullet \mapsto H^i(K^\bullet)$ se prolonge en des foncteurs $H^i :
K(\mathcal{A}) \to \mathcal{A}$ et $H^i : D(\mathcal{A}) \to \mathcal{A}$.

\begin{lemma}
\label{lemma-when-in-bounded}
Un objet $E$ de $D(\mathcal{A})$ appartient à $D^+(\mathcal{A})$ si et seulement
si $H^i(E) =0 $ pour tout $i \ll 0$. Des énoncés analogues valent pour $D^-$ et
$D^b$.
\end{lemma}

\begin{proof}
Indication : utiliser les foncteurs de troncation. Voir
Catégories dérivées, lemme \ref{derived-lemma-complex-cohomology-bounded}.
\end{proof}

\begin{lemma}
\label{lemma-derived-categories}
Morphismes entre objets de la catégorie dérivée.
\begin{enumerate}
\item
Soit $I^\bullet \in \text{Comp}^+(\mathcal{A})$ tel que $I^n$ soit injectif pour tout
$n \in \mathbf{Z}$. Alors
$$
\Hom_{D(\mathcal{A})}(K^\bullet, I^\bullet)
=
\Hom_{K(\mathcal{A})}(K^\bullet, I^\bullet).
$$
\item
Soit $P^\bullet \in \text{Comp}^-(\mathcal{A})$ tel que $P^n$ soit projectif pour tout
$n \in \mathbf{Z}$. Alors
$$
\Hom_{D(\mathcal{A})}(P^\bullet, K^\bullet)
=
\Hom_{K(\mathcal{A})}(P^\bullet, K^\bullet).
$$
\item
Si $\mathcal{A}$ possède assez d'injectifs et si $\mathcal{I} \subset \mathcal{A}$
est la sous-catégorie additive des objets injectifs, alors
$
D^+(\mathcal{A})\cong K^+(\mathcal{I})
$
(en tant que catégories triangulées).
\item
Si $\mathcal{A}$ possède assez de projectifs et si $\mathcal{P} \subset \mathcal{A}$
est la sous-catégorie additive des objets projectifs, alors
$
D^-(\mathcal{A}) \cong K^-(\mathcal{P}).
$
\end{enumerate}
\end{lemma}

\begin{proof}
Démonstration omise.
\end{proof}

\begin{definition}
\label{definition-derived-functor}
Soit $F: \mathcal{A} \to \mathcal{B}$ un foncteur exact à gauche et supposons que
$\mathcal{A}$ possède assez d'injectifs. Nous définissons le {\it foncteur dérivé total
à droite de $F$} comme le foncteur $RF: D^+(\mathcal{A}) \to D^+(\mathcal{B})$
s'insérant dans le diagramme
$$
\xymatrix{
D^+(\mathcal{A}) \ar[r]^{RF} & D^+(\mathcal{B}) \\
K^+(\mathcal I) \ar[u] \ar[r]^F & K^+(\mathcal{B}). \ar[u]
}
$$
Cela est possible puisque la flèche verticale de gauche est inversible d'après le lemme
précédent. De même, soit $G: \mathcal{A} \to \mathcal{B}$ un foncteur exact à droite
et supposons que $\mathcal{A}$ possède assez de projectifs. Nous définissons le
{\it foncteur dérivé total à gauche de $G$} comme le foncteur $LG: D^-(\mathcal{A})
\to D^-(\mathcal{B})$ s'insérant dans le diagramme
$$
\xymatrix{
D^-(\mathcal{A}) \ar[r]^{LG} & D^-(\mathcal{B}) \\
K^-(\mathcal{P}) \ar[u] \ar[r]^G & K^-(\mathcal{B}). \ar[u]
}
$$
Cela est possible puisque la flèche verticale de gauche est inversible d'après le lemme
précédent.
\end{definition}

\begin{remark}
\label{remark-cohomology-of-derived-functor}
Dans ces cas, on a $R^iF(K^\bullet) = H^i(RF(K^\bullet))$, où
le membre de gauche est défini comme la $i$-ième homologie du complexe
$F(K^\bullet)$.
\end{remark}




\section{Catégorie dérivée filtrée}
\label{section-filtered-derived-category}

\noindent
Il s'avère que nous devons tout recommencer et construire également la catégorie
dérivée filtrée.

\begin{definition}
\label{definition-filtered}
Soit $\mathcal{A}$ une catégorie abélienne.
\begin{enumerate}
\item Soit $\text{Fil}(\mathcal{A})$ la catégorie des objets filtrés
$(A, F)$ de $\mathcal{A}$, où $F$ est une filtration de la forme
$$
A \supset \ldots \supset F^n A \supset F^{n+1}A \supset \ldots
\supset 0.
$$
Il s'agit d'une catégorie additive.
\item Nous notons $\text{Fil}^f(\mathcal{A})$ la sous-catégorie pleine
de $\text{Fil}(\mathcal{A})$ dont les objets $(A, F)$ ont une filtration
finie. C'est encore une catégorie additive.
\item Un objet $I \in \text{Fil}^f(\mathcal{A})$ est dit
{\it injectif filtré} (respectivement {\it projectif filtré}) si
$\text{gr}^p(I) = \text{gr}_F^p(I) = F^pI/F^{p+1}I$ est injectif
(respectivement projectif) dans $\mathcal{A}$ pour tout $p$.
\item La catégorie des complexes
$\text{Comp}(\text{Fil}^f(\mathcal{A})) \supset
\text{Comp}^+(\text{Fil}^f(\mathcal{A}))$
et sa catégorie homotopique
$K(\text{Fil}^f(\mathcal{A})) \supset K^+(\text{Fil}^f(\mathcal A))$
sont définies comme précédemment.
\item Un morphisme $\alpha : K^\bullet \to L^\bullet$ de complexes de
$\text{Comp}(\text{Fil}^f(\mathcal{A}))$ est appelé un
{\it quasi-isomorphisme filtré} si
$$
\text{gr}^p(\alpha): \text{gr}^p(K^\bullet) \to \text{gr}^p(L^\bullet)
$$
est un quasi-isomorphisme pour tout $p \in \mathbf{Z}$.
\item Nous définissons $DF(\mathcal{A})$ (respectivement $DF^+(\mathcal{A})$)
en inversant les quasi-isomorphismes filtrés dans
$K(\text{Fil}^f(\mathcal{A}))$ (respectivement $K^+(\text{Fil}^f(\mathcal{A}))$).
\end{enumerate}
\end{definition}

\begin{lemma}
\label{lemma-filtered-derived-category}
Si $\mathcal{A}$ possède assez d'injectifs, alors $DF^+(\mathcal{A}) \cong
K^+(\mathcal{I})$, où $\mathcal{I}$ est la sous-catégorie additive pleine de
$\text{Fil}^f(\mathcal{A})$ formée des objets injectifs filtrés.
De même, si $\mathcal{A}$ possède assez de projectifs, alors $DF^-(\mathcal{A})
\cong K^-(\mathcal{P})$, où $\mathcal P$ est la sous-catégorie additive pleine de
$\text{Fil}^f(\mathcal{A})$ formée des objets projectifs filtrés.
\end{lemma}

\begin{proof}
Démonstration omise.
\end{proof}





\section{Foncteurs dérivés filtrés}
\label{section-filtered-derived-functors}

\noindent
Il y a ensuite les foncteurs dérivés filtrés.

\begin{definition}
\label{definition-filtered-derived-functors}
Soit $T: \mathcal{A} \to \mathcal{B}$ un foncteur exact à gauche et supposons que
$\mathcal{A}$ possède assez d'injectifs. Définissons $RT: DF^+(\mathcal{A}) \to D
F^+(\mathcal{B})$ de sorte qu'il s'insère dans le diagramme
$$
\xymatrix{
DF^+(\mathcal{A}) \ar[r]^{RT} & DF^+(\mathcal{B}) \\
K^+(\mathcal{I}) \ar[u] \ar[r]^{T \quad} & K^+(\text{Fil}^f(\mathcal{B})).
\ar[u]}
$$
Ce foncteur est bien défini d'après le lemme précédent. Soit $G: \mathcal{A} \to
\mathcal{B}$ un foncteur exact à droite et supposons que $\mathcal{A}$ possède assez de
projectifs. Définissons $LG: DF^-(\mathcal{A}) \to DF^-(\mathcal{B})$ de sorte qu'il s'insère dans
le diagramme
$$
\xymatrix{
DF^-(\mathcal{A}) \ar[r]^{LG} & DF^-(\mathcal{B}) \\
K^-(\mathcal{P}) \ar[u] \ar[r]^{G \quad} & K^-(\text{Fil}^f(\mathcal{B})).
\ar[u]}
$$
Là encore, ce foncteur est bien défini d'après le lemme précédent.
Les foncteurs $RT$, respectivement\ $LG$, sont appelés le {\it foncteur dérivé
filtré} de $T$, respectivement de $G$.
\end{definition}

\begin{proposition}
\label{proposition-compare-filtered-graded}
Dans la situation ci-dessus, on a
$$
\text{gr}^p \circ RT = RT \circ \text{gr}^p
$$
où le $RT$ de gauche est le foncteur dérivé filtré, tandis que celui de
droite est le foncteur dérivé total. Autrement dit, on a un diagramme commutatif
$$
\xymatrix{
DF^+(\mathcal{A}) \ar[r]^{RT} \ar[d]_{\text{gr}^p} & DF^+(\mathcal{B})
\ar[d]^{\text{gr}^p}\\
D^+(\mathcal{A}) \ar[r]^{RT} & D^+(\mathcal{B}).}
$$
\end{proposition}

\begin{proof}
Démonstration omise.
\end{proof}

\noindent
Pour $K^\bullet \in DF^+(\mathcal{B})$, on obtient une suite spectrale
$$
E_1^{p, q} = H^{p+q}(\text{gr}^p K^\bullet) \Rightarrow H^{p+q}(\text{oubli de la
filtration}(K^\bullet)).
$$






\section{Application des complexes filtrés}
\label{section-applications-filtered}

\noindent
Soit $\mathcal{A}$ une catégorie abélienne possédant assez d'injectifs, et soit
$0 \to L \to M \to N \to 0$ une suite exacte courte dans $\mathcal{A}$.
Considérons $\widetilde M \in \text{Fil}^f(\mathcal{A})$ obtenu en munissant $M$ de la
filtration définie par
$$
F^1M = L, \ F^nM = M
\text{ pour }n \leq 0\text{, et }F^nM = 0\text{ pour }n \geq 2.
$$
Par définition, on a
$$
\text{oubli de la filtration}(\widetilde M) = M, \quad
\text{gr}^0(\widetilde M) = N, \quad
\text{gr}^1(\widetilde M) = L
$$
et $\text{gr}^n(\widetilde M) = 0$ pour tout autre $n \neq 0, 1$. Soit $T:
\mathcal{A} \to \mathcal{B}$ un foncteur exact à gauche. Supposons que $\mathcal{A}$
possède assez d'injectifs. Alors $RT(\widetilde M) \in DF^+(\mathcal{B})$ est un
complexe filtré vérifiant
$$
\text{gr}^p(RT(\widetilde M))
\stackrel{\text{qis}}{=}
\left\{
\begin{matrix}
0 & \text{si} & p \neq 0, 1, \\
RT(N) & \text{si} & p = 0, \\
RT(L) & \text{si} & p = 1.
\end{matrix}
\right.
$$
et $\text{oubli de la filtration}(RT(\widetilde M))\stackrel{\text{qis}}{ = } RT(M)$. La
suite spectrale appliquée à $RT(\widetilde M)$ donne
$$
E_1^{p, q} = R^{p+q}T(\text{gr}^p(\widetilde M)) \Rightarrow
R^{p+q}T(\text{oubli de la filtration}(\widetilde M)).
$$
En développant la suite spectrale, on obtient la suite exacte longue
$$
\xymatrix{
0 \ar[r] & T(L) \ar[r] & T(M) \ar[r] & T(N) \ar@(rd, ul)[rdllllr] \\
& R^1T(L) \ar[r] & R^1T(M) \ar[r] & \ldots
}
$$
Nous l'utiliserons de la manière suivante. Soit $X/k$ un schéma de type fini. Soit
$\mathcal{F}$ un $\mathbf{Z}/\ell^n \mathbf{Z}$-module constructible plat.
Nous voulons alors montrer que la trace
$$
\text{Tr}( \pi_X^\ast | R\Gamma_c(X_{\bar k}, \mathcal{F})) \in
\mathbf{Z}/\ell^n \mathbf{Z}
$$
est additive sur les suites exactes courtes. Pour le voir, il ne suffira pas de
travailler avec $R\Gamma_c(X_{\bar k}, -) \in D^+(\mathbf{Z}/\ell^n \mathbf{Z})$, mais
il faudra utiliser la catégorie dérivée filtrée.







%10.29.09
\section{Complexes parfaits}
\label{section-perfect}

\noindent
Soit $\Lambda$ un anneau (éventuellement non commutatif), $\text{Mod}_{\Lambda}$ la
catégorie des $\Lambda$-modules à gauche, $K(\Lambda) = K(\text{Mod}_\Lambda)$ sa
catégorie homotopique, et $D(\Lambda)= D(\text{Mod}_\Lambda)$ la catégorie dérivée
correspondante.

\begin{definition}
\label{definition-perfect}
Nous notons $K_{perf}(\Lambda)$ la catégorie dont les objets sont les complexes bornés
de $\Lambda$-modules projectifs de type fini, et dont les morphismes sont les
morphismes de complexes à homotopie près. Le foncteur $K_{perf}(\Lambda)\to
D(\Lambda)$ est pleinement fidèle (Catégories dérivées, lemme
\ref{derived-lemma-morphisms-from-projective-complex}).
Notons $D_{perf}(\Lambda)$ son image essentielle.
Un objet de $D(\Lambda)$ est dit {\it parfait} s'il appartient à
$D_{perf}(\Lambda)$.
\end{definition}

\begin{proposition}
\label{proposition-trace-well-defined}
Soient $K\in D_{perf}(\Lambda)$ et $f\in \text{End}_{D(\Lambda)}(K)$. Alors la
trace $\text{Tr}(f)\in \Lambda^\natural$ est bien définie.
\end{proposition}

\begin{proof}
Nous utiliserons le lemme de Catégories dérivées
\ref{derived-lemma-morphisms-from-projective-complex}
sans plus le mentionner dans cette démonstration.
Soit $P^\bullet$ un complexe borné de $\Lambda$-modules projectifs de type fini
et soit $\alpha : P^\bullet \to K$ un isomorphisme dans $D(\Lambda)$. Alors
$\alpha^{-1}\circ f\circ \alpha$ correspond à un morphisme de complexes
$f^\bullet : P^\bullet \to P^\bullet$ bien défini à homotopie près.
Posons
$$
\text{Tr}(f) = \sum_i (-1)^i \text{Tr}(f^i : P^i \to P^i) \in \Lambda^\natural.
$$
Une fois $P^\bullet$ et $\alpha$ fixés, cette définition ne dépend pas du choix de
$f^\bullet$. En effet, tout autre choix est de la forme
$\tilde{f}^\bullet = f^\bullet + dh +hd$ pour un certain
$h^i : P^i \to P^{i-1}(i\in \mathbf{Z})$. Mais
\begin{eqnarray*}
\text{Tr}(dh) & = & \sum_i (-1)^i \text{Tr}(P^i\xrightarrow{dh} P^i) \\
& = & \sum_i (-1)^i \text{Tr}(P^{i-1}\xrightarrow{hd} P^{i-1}) \\
& = & -\sum_i (-1)^{i-1}\text{Tr}(P^{i-1}\xrightarrow{hd} P^{i-1}) \\
& = & - \text{Tr}(hd)
\end{eqnarray*}
donc $\sum_i (-1)^i \text{Tr} ((dh+hd)|_{P^i})=0$.
En outre, cette définition ne dépend pas du choix de $(P^\bullet , \alpha)$ :
supposons que $(Q^\bullet, \beta)$ soit un autre choix. Les composées
$$
Q^\bullet \xrightarrow{\beta} K \xrightarrow{\alpha^{-1}} P^\bullet
\quad\text{et}\quad
P^\bullet \xrightarrow{\alpha} K \xrightarrow{\beta^{-1}} Q^\bullet
$$
sont représentables par des morphismes de complexes $\gamma_1^\bullet$ et
$\gamma_2^\bullet$ respectivement, tels que $\gamma_1^\bullet \circ
\gamma_2^\bullet$ soit homotope à l'identité. Ainsi, le morphisme de complexes
$\gamma_2^\bullet\circ f^\bullet\circ \gamma_1^\bullet : Q^\bullet\to Q^\bullet$
représente le morphisme $\beta^{-1}\circ f\circ\beta$ dans $D(\Lambda)$. Or
\begin{eqnarray*}
\text{Tr}(\gamma_2^\bullet\circ f^\bullet\circ\gamma_1^\bullet|_{Q^\bullet}) &
= & \text{Tr}(\gamma_1^\bullet \circ\gamma_2^\bullet \circ
f^\bullet|_{P^\bullet})\\
& = & \text{Tr}(f^\bullet|_{P^\bullet})
\end{eqnarray*}
car $\gamma_1^\bullet \circ \gamma_2^\bullet$ est homotope à
l'identité et d'après l'indépendance du choix de $f^\bullet$ établie ci-dessus.
\end{proof}




\section{Filtrations et complexes parfaits}
\label{section-filtrations-perfect}

\noindent
Présentons maintenant une version filtrée de la catégorie des complexes parfaits. Un
objet $(M, F)$ de $\text{Fil}^f(\text{Mod}_\Lambda)$ est dit {\it projectif filtré
de type fini} si, pour tout $p$, $\text{gr}^p_F (M)$ est de type fini et
projectif. Nous considérons alors la catégorie homotopique
$KF_{\text{perf}}(\Lambda)$ des complexes bornés d'objets projectifs filtrés
de type fini de $\text{Fil}^f(\text{Mod}_\Lambda)$. On a un diagramme de
catégories
$$
\begin{matrix}
KF(\Lambda) & \supset & KF_{\text{perf}}(\Lambda)\\
\downarrow & & \downarrow\\
DF(\Lambda) & \supset & DF_{\text{perf}}(\Lambda)
\end{matrix}
$$
où le foncteur vertical de droite est pleinement fidèle et où la catégorie
$DF_{\text{perf}}(\Lambda)$ est son image essentielle, comme précédemment.

\begin{lemma}[Additivité]
\label{lemma-additivity}
Soient $K\in DF_{\text{perf}}(\Lambda)$ et $f\in
\text{End}_{DF}(K)$. Alors
$$
\text{Tr}(f|_K) =
\sum\nolimits_{p\in \mathbf{Z}} \text{Tr}(f|_{\text{gr}^p K}).
$$
\end{lemma}

\begin{proof}
D'après la proposition \ref{proposition-trace-well-defined}, on peut supposer disposer
d'un complexe borné
$P^\bullet$ d'objets projectifs filtrés de type fini de
$\text{Fil}^f(\text{Mod}_\Lambda)$ et d'un morphisme $f^\bullet : P^\bullet\to
P^\bullet$ dans $\text{Comp}(\text{Fil}^f(\text{Mod}_\Lambda))$. Le lemme
découle donc du résultat suivant, dont la démonstration est laissée au lecteur.
\end{proof}

\begin{lemma}
\label{lemma-additive-filtered-finite-projective}
Soit $P \in \text{Fil}^f(\text{Mod}_\Lambda)$ un objet projectif filtré de type fini, et
soit $f : P \to P$ un endomorphisme dans $\text{Fil}^f(\text{Mod}_\Lambda)$. Alors
$$
\text{Tr}(f|_P) =
\sum\nolimits_p \text{Tr}(f|_{\text{gr}^p(P)}).
$$
\end{lemma}

\begin{proof}
Démonstration omise.
\end{proof}







\section{Caractérisation des objets parfaits}
\label{section-characterizing-perfect}

\noindent
Pour le cas commutatif, voir
Compléments sur l'algèbre, sections
\ref{more-algebra-section-pseudo-coherent},
\ref{more-algebra-section-tor}, et
\ref{more-algebra-section-perfect}.

\begin{definition}
\label{definition-finite-tor-dimension}
Soit $\Lambda$ un anneau (éventuellement non commutatif).
Un objet $K\in D(\Lambda)$ est de {\it dimension de $\text{Tor}$ finie}
s'il existe $a, b \in \mathbf{Z}$ tels que, pour tout
$\Lambda$-module à droite $N$, on ait
$H^i(N \otimes_{\Lambda}^\mathbf{L} K) = 0$ pour tout
$i \not \in [a, b]$.
\end{definition}

\noindent
Cela signifie en particulier que $K \in D^b(\Lambda)$, comme on le voit en prenant
$N = \Lambda$.

\begin{lemma}
\label{lemma-characterize-perfect}
Soient $\Lambda$ un anneau noethérien à gauche et $K\in D(\Lambda)$. Alors $K$ est
parfait si et seulement si les deux conditions suivantes sont satisfaites :
\begin{enumerate}
\item
$K$ est de dimension de $\text{Tor}$ finie, et
\item
pour tout $i \in \mathbf{Z}$, $H^i(K)$ est un $\Lambda$-module de type fini.
\end{enumerate}
\end{lemma}

\begin{proof}
Voir Compléments sur l'algèbre, lemme \ref{more-algebra-lemma-perfect},
pour la démonstration dans le cas commutatif.
\end{proof}

\noindent
Le lecteur est vivement encouragé à essayer de démontrer cet énoncé. La démonstration repose sur le
fait qu'un module à gauche de présentation finie sur un anneau est plat si et seulement
s'il est projectif.

\begin{remark}
\label{remark-variant}
Une variante de ce lemme consiste à considérer un schéma noethérien $X$
et la catégorie $D_{perf}(\mathcal{O}_X)$ des complexes qui sont localement
quasi-isomorphes à un complexe fini de modules localement libres de type fini sur
$\mathcal{O}_X$. Les objets $K$ de $D_{perf}(\mathcal{O}_X)$
se caractérisent par la cohérence de leurs faisceaux de cohomologie et par une
dimension de Tor bornée.
\end{remark}








\section{Cohomologie des complexes raisonnables}
\label{section-cohomology-ctf}

\noindent
Ce qui suit est un cas particulier d'un résultat plus général concernant la
cohomologie à support propre des objets de $D_{ctf}(X, \Lambda)$.

\begin{proposition}
\label{proposition-projective-curve-constructible-cohomology}
Soient $X$ une courbe projective sur un corps $k$, $\Lambda$ un anneau fini et
$K\in D_{ctf}(X, \Lambda)$. Alors $R\Gamma(X_{\bar k}, K)\in
D_{perf}(\Lambda)$.
\end{proposition}

\begin{proof}[Esquisse de démonstration.]
La première étape consiste à montrer :
\begin{enumerate}
\item[(1)]
{\it La cohomologie de $R\Gamma(X_{\bar k}, K)$ est bornée.}
\end{enumerate}
Considérons la suite spectrale
$$
H^i(X_{\bar k}, \underline H^j(K))
\Rightarrow
H^{i+j} (R\Gamma(X_{\bar k}, K)).
$$
Comme $K$ est borné et que $\Lambda$ est fini, les faisceaux $\underline H^j(K)$
sont de torsion. De plus, $X_{\bar k}$ est de dimension cohomologique finie, donc le
membre de gauche n'est non nul que pour un nombre fini de $i$ et de $j$. Il en va donc
de même du membre de droite.
\begin{enumerate}
\item[(2)]
{\it Les groupes de cohomologie $H^{i+j} (R\Gamma(X_{\bar k}, K))$ sont finis.}
\end{enumerate}
Comme les faisceaux $\underline H^j(K)$ sont constructibles, les groupes
$H^i(X_{\bar k}, \underline H^j(K))$ sont finis
(Cohomologie \'etale, section \ref{etale-cohomology-section-vanishing-torsion})
et l'assertion découle de nouveau de la suite spectrale.
\begin{enumerate}
\item[(3)]
{\it $R\Gamma(X_{\bar k}, K)$ est de dimension de $\text{Tor}$ finie.}
\end{enumerate}
Soit $N$ un $\Lambda$-module à droite (en fait, comme $\Lambda$ est fini, il
suffit de supposer que $N$ est fini). Par la formule de projection (changement de
module),
$$
N \otimes^\mathbf{L}_\Lambda R \Gamma(X_{\bar k}, K) = R\Gamma(X_{\bar k},
N \otimes^\mathbf{L}_\Lambda K).
$$
Par conséquent,
$$
H^i (N \otimes^\mathbf{L}_\Lambda R\Gamma(X_{\bar k}, K)) = H^i(R\Gamma(X_{\bar
k}, N \otimes_{\Lambda}^\mathbf{L} K)).
$$
Considérons maintenant la suite spectrale
$$
H^i (X_{\bar k}, \underline H^j (N \otimes_{\Lambda}^\mathbf{L} K))
\Rightarrow
H^{i+j}(R\Gamma(X_{\bar k}, N \otimes_{\Lambda}^\mathbf{L} K)).
$$
Comme $K$ est de dimension de $\text{Tor}$ finie, $\underline H^j
(N \otimes_{\Lambda}^\mathbf{L} K)$ s'annule universellement pour $j$ assez petit,
et le membre de gauche s'annule dès que $i < 0$. Par conséquent, $R\Gamma(X_{\bar
k}, K)$ est de dimension de $\text{Tor}$ finie, comme annoncé. C'est donc un
complexe parfait d'après le lemme \ref{lemma-characterize-perfect}.
\end{proof}





\section{Nombres de Lefschetz}
\label{section-lefschetz-numbers}

\noindent
Le fait que la cohomologie totale d'un complexe constructible de dimension de Tor
finie soit un complexe parfait est la raison technique essentielle pour laquelle la cohomologie
se comporte bien et permet de définir rigoureusement les traces qui apparaissent dans la
formule des traces.

\begin{definition}
\label{definition-global-lefschetz-number}
Soient $\Lambda$ un anneau fini, $X$ une courbe projective sur un corps fini $k$
et $K \in D_{ctf}(X, \Lambda)$ (par exemple $K = \underline\Lambda$).
Il existe un morphisme canonique $c_K : \pi_X^{-1}K \to K$, et son changement de base
$c_K|_{X_{\bar k}}$ induit une action, notée $\pi_X^*$, sur le complexe parfait
$R\Gamma(X_{\bar k}, K|_{X_{\bar k}})$. Le
{\it nombre global de Lefschetz} de $K$ est la trace
$\text{Tr}(\pi_X^* |_{R\Gamma(X_{\bar k}, K)})$ de cette action.
C'est un élément de $\Lambda^\natural$.
\end{definition}

\begin{definition}
\label{definition-local-lefschetz-number}
Soient $\Lambda, X, k, K$ comme dans la
définition \ref{definition-global-lefschetz-number}.
Comme $K\in D_{ctf}(X, \Lambda)$, pour tout point géométrique $\bar x$ de $X$,
le complexe $K_{\bar x}$ est un complexe parfait (dans $D_{perf}(\Lambda)$). Comme nous
l'avons vu à la section \ref{section-frobenii}, le Frobenius $\pi_X$ agit sur
$K_{\bar x}$. Le {\it nombre local de Lefschetz} de $K$ est la somme
$$
\sum\nolimits_{x\in X(k)} \text{Tr}(\pi_x |_{K_{\overline{x}}})
$$
qui est encore un élément de $\Lambda^\natural$.
\end{definition}

\noindent
Nous pouvons enfin formuler précisément la formule des traces.

\begin{theorem}[Formule des traces de Lefschetz]
\label{theorem-trace}
Soient $X$ une courbe projective sur un corps fini $k$, $\Lambda$ un anneau fini
et $K \in D_{ctf}(X, \Lambda)$. Alors les nombres de Lefschetz global et local
de $K$ sont égaux, c'est-à-dire
\begin{equation}
\label{equation-trace-formula}
\text{Tr}(\pi^*_X |_{R\Gamma(X_{\bar k}, K)})
=
\sum\nolimits_{x\in X(k)} \text{Tr}(\pi_X |_{K_{\bar x}})
\end{equation}
dans $\Lambda^\natural$.
\end{theorem}

\begin{proof}
Voir la discussion ci-dessous.
\end{proof}

%11.5.09
\noindent
Nous utiliserons la formule des traces plutôt que de la démontrer. Néanmoins, nous
donnerons de nombreux détails de la démonstration du théorème telle qu'elle est donnée dans
\cite{SGA4.5} (certaines des choses qui n'y sont pas expliquées de manière satisfaisante
sont énumérées à la section \ref{section-list-skipped}).

\medskip\noindent
Nous n'avons énoncé la formule que pour les courbes et, en un sens assez
faible, elle est une conséquence du résultat suivant.

\begin{theorem}[Weil]
\label{theorem-weil-trace-formula}
Soient $C$ une courbe projective non singulière sur un corps algébriquement clos
$k$, et $\varphi : C \to C$ un $k$-endomorphisme de $C$ distinct de
l'identité. Posons $V(\varphi) = \Delta_C \cdot \Gamma_\varphi$, où $\Delta_C$ est
la diagonale, $\Gamma_\varphi$ est le graphe de $\varphi$, et où le nombre d'intersection
est calculé sur $C \times C$. Soit $J = \underline{\Picardfunctor}^0_{C/k}$
la jacobienne de $C$, et notons $\varphi^* : J \to J$ l'action induite par
$\varphi$ par image réciproque. Alors
$$
V(\varphi) = 1 - \text{Tr}_J(\varphi^*) + \deg \varphi.
$$
\end{theorem}

\begin{proof}
Le nombre $V(\varphi)$ est le nombre de points fixes de $\varphi$ comptés avec multiplicité ; il est égal
à
$$
V(\varphi) =
\sum\nolimits_{c \in |C| : \varphi(c) = c} m_{\text{Fix}(\varphi)} (c)
$$
où $m_{\text{Fix}(\varphi)} (c)$ est la multiplicité de $c$ comme point fixe
de $\varphi$, c'est-à-dire l'ordre d'annulation de l'image d'une uniformisante locale
par $\varphi - \text{id}_C$. On trouvera des démonstrations de ce théorème dans
\cite{Lang} et \cite{Weil}.
\end{proof}

\begin{example}
\label{example-elliptic-curve}
Soit $C = E$ une courbe elliptique et soit $\varphi = [n]$ la multiplication par $n$.
Alors $\varphi^* = \varphi^t$ est la multiplication par $n$ sur la jacobienne, de sorte qu'elle
a pour trace $2n$ et pour degré $n^2$. D'autre part, le schéma des points fixes de
$\varphi$ est formé des points $p \in E$ tels que $n p = p$ ; c'est donc le schéma de
$(n-1)$-torsion, qui est de degré $(n-1)^2$. Le théorème s'écrit donc
$$
(n-1)^2 = 1 - 2n + n^2.
$$
\end{example}

\noindent
{\bf Jacobiennes.}
Nous examinons maintenant, sans démonstration, la correspondance entre une courbe et sa
jacobienne qui intervient dans la démonstration de Weil. Soit $C$ une courbe projective non singulière
sur un corps algébriquement clos $k$ et choisissons un point de base $c_0 \in
C(k)$. Notons $A^1(C \times C)$ (ou $\Pic(C \times C)$, ou
$\text{CaCl}(C \times C)$) le groupe abélien des diviseurs de codimension 1 de
$C \times C$. Alors
$$
A^1(C \times C) = \text{pr}_1^* (A^1(C)) \oplus \text{pr}_2^* (A^1(C)) \oplus R
$$
où
$$
R = \{ Z \in A^1(C \times C) \ |
\ Z|_{C \times \{c_0\}} \sim_\text{rat} 0
\text{ et }
Z|_{\{c_0\} \times C} \sim_\text{rat} 0 \}.
$$
Autrement dit,
$R$ est le sous-groupe des fibrés en droites dont les restrictions aux
deux sous-schémas indiqués sont triviales. Il existe alors un isomorphisme canonique de groupes abéliens $R
\cong \text{End}(J)$ qui envoie un diviseur $Z$ de $R$ sur l'endomorphisme
$$
\begin{matrix}
J & \to & J \\
\left[ \mathcal{O}_C(D) \right] & \mapsto & (\text{pr}_1 |_Z)_* (\text{pr}_2
|_Z)^* (D).
\end{matrix}
$$
La correspondance susmentionnée est la suivante. Notons $\sigma$
l'automorphisme de $C \times C$ qui permute les facteurs.
$$
\begin{matrix}
\hline & \\
\text{End}(J) & R \\
& \\
\hline & \\
\text{composition de }\alpha, \beta &
{\text{pr}_{13}}_* ({\text{pr}_{12}}^*(\alpha) \circ {\text{pr}_{23}}^*(\beta))
\\
& \\
\text{id}_J &
\Delta_C - \{c_0\} \times C - C \times \{c_0\} \\
& \\
\varphi^* &
\Gamma_\varphi - C \times \{\varphi(c_0)\}
- \sum_{\varphi(c) = c_0} \{c\} \times C \\
& \\
{
\begin{matrix}
\text{la forme trace} \\
\alpha, \beta \mapsto \text{Tr}(\alpha \beta)
\end{matrix}
}
&
\alpha, \beta \mapsto - \int_{C \times C} \alpha . \sigma^*\beta
\\
& \\
{
\begin{matrix}
\text{l'involution de Rosati} \\
\alpha \mapsto \alpha^\dagger
\end{matrix}
}
&
\alpha \mapsto \sigma^*\alpha
\\
& \\
{
\begin{matrix}
\text{positivité de Rosati} \\
\text{Tr}(\alpha\alpha^\dagger) > 0
\end{matrix}
}
&
{
\begin{matrix}
\text{théorème de l'indice de Hodge sur }C \times C \\
- \int_{C \times C} \alpha \sigma^*\alpha > 0.
\end{matrix}
}
\\
& \\
\hline 
\end{matrix}
$$
En fait, au vu de la formule de Künneth, le sous-groupe $R$ correspond aux classes de Hodge de type
$1, 1$ dans $H^1(C)\otimes H^1(C)$.

\medskip\noindent
{\bf Démonstration de Weil.} Cette correspondance permet de démontrer la formule
des traces. Nous avons
\begin{eqnarray*}
V(\varphi) & = & \int_{C \times C} \Gamma_\varphi.\Delta \\
& = & \int_{C \times C} \Gamma_\varphi. \left(\Delta_C - \{c_0\} \times C - C
\times \{c_0\}\right) + \int_{C \times C} \Gamma_\varphi. \left(\{c_0\} \times C
+ C \times \{c_0\}\right).
\end{eqnarray*}
D'une part,
$$
\int_{C \times C} \Gamma_\varphi. \left(\{c_0\} \times C + C \times
\{c_0\}\right)
=
1 + \deg \varphi
$$
et, d'autre part, puisque $R$ est l'orthogonal du diviseur ample
$\{c_0\} \times C + C \times \{c_0\}$,
\begin{eqnarray*}
& &
\int_{C \times C} \Gamma_\varphi. \left(\Delta_C - \{c_0\} \times C - C \times
\{c_0\}\right) \\
& = &
\int_{C \times C} \left(\Gamma_\varphi - C \times \{\varphi(c_0)\} -
\sum_{\varphi(c) = c_0} \{c\} \times C \right). \left(\Delta_C - \{c_0\} \times
C - C \times \{c_0\}\right) \\
& = & - \text{Tr}_J (\varphi^* \circ \text{id}_J).
\end{eqnarray*}
En récapitulant, nous obtenons
$$
V(\varphi) = 1 - \text{Tr}_J (\varphi^*) + \deg \varphi
$$
ce qui donne la formule des traces.

\begin{lemma}
\label{lemma-weil-mod}
Considérons la situation du
théorème \ref{theorem-weil-trace-formula}
et soit $\ell$ un nombre premier inversible dans $k$. Alors
$$
\sum\nolimits_{i = 0}^2
(-1)^i
\text{Tr}(\varphi^* |_{H^i (C, \underline{\mathbf{Z}/\ell^n \mathbf{Z}})})
=
V(\varphi) \mod \ell^n.
$$
\end{lemma}

\begin{proof}[Esquisse de démonstration]
Remarquons d'abord que l'énoncé a un sens, car les groupes $H^i(C,
\underline{\mathbf{Z}/\ell^n \mathbf{Z}})$ sont des modules libres sur $\mathbf{Z}/\ell^n
\mathbf{Z}$ pour tout $i$. La trace de $\varphi^*$ sur la
cohomologie de degré zéro vaut 1. Le choix d'une racine primitive $\ell^n$-ième de l'unité dans $k$
fournit un isomorphisme
$$
H^i(C, \underline{\mathbf{Z}/\ell^n \mathbf{Z}}) \cong H^i(C, \mu_{\ell^n})
$$
de manière compatible avec l'action induite par l'endomorphisme considéré. D'autre part,
$H^1(C, \mu_{\ell^n}) = J[\ell^n]$. Par conséquent,
\begin{eqnarray*}
\text{Tr}(\varphi^* |_{H^1 (C, \underline{\mathbf{Z}/\ell^n \mathbf{Z}})})) & =
& \text{Tr}_J (\varphi^*) \mod \ell^n \\
& = & \text{Tr}_{\mathbf{Z}/\ell^n \mathbf{Z}} (\varphi^* : J[\ell^n] \to
J[\ell^n]).
\end{eqnarray*}
En outre, $H^2(C, \mu_{\ell^n}) = \Pic(C)/\ell^n\Pic(C) \cong
\mathbf{Z}/\ell^n \mathbf{Z}$, où $\varphi^*$ agit par multiplication par $\deg
\varphi$. Par suite,
$$
\text{Tr} (\varphi^*|_{H^2 (C, \underline{\mathbf{Z}/\ell^n \mathbf{Z}})}) =
\deg \varphi.
$$
Ainsi,
$$
\sum_{i = 0}^2 (-1)^i
\text{Tr}(\varphi^* |_{H^i (C, \underline{\mathbf{Z}/\ell^n \mathbf{Z}})}) =
1 - \text{Tr}_J(\varphi^*) + \deg \varphi \mod \ell^n
$$
et le lemme découle du théorème \ref{theorem-weil-trace-formula}.
\end{proof}

\noindent
Une autre manière de démontrer ce lemme consiste à montrer que
$$
X \mapsto H^* (X, \mathbf{Q}_\ell) =
\mathbf{Q}_\ell \otimes
\lim_n H^*(X, \mathbf{Z}/\ell^n\mathbf{Z})
$$
définit une théorie cohomologique de Weil sur les variétés projectives lisses sur $k$. Alors
la formule des traces
$$
V(\varphi) = \sum_{i = 0}^2 (-1)^i
\text{Tr}(\varphi^* |_{H^i(C, \mathbf{Q}_\ell)})
$$
est une conséquence formelle des axiomes (il s'agit d'un exercice d'algèbre linéaire, la
démonstration étant la même que dans le cas topologique).




%11.10.09
\section{Préliminaires et sorites}
\label{section-preliminaries}

\noindent
Notations :
Fixons les notations de cette section. On note $A$ un anneau commutatif,
et $\Lambda$ un anneau (éventuellement non commutatif) muni d'un morphisme d'anneaux $A\to \Lambda$ dont
l'image est contenue dans le centre de $\Lambda$. Soient $G$ un groupe fini et $\Gamma$ une
{\it extension de monoïdes de noyau $G$ et de quotient $\mathbf{N}$}, c'est-à-dire qu'il existe une
suite exacte
$$
1\to G\to \tilde\Gamma\to \mathbf{Z}\to 1
$$
et $\Gamma$ est constitué des éléments de $\tilde\Gamma$ dont l'image est
positive ou nulle. Enfin, soit $P$ un $A[\Gamma]$-module qui est de type fini et
projectif en tant que $A[G]$-module, et soit $M$ un $\Lambda[\Gamma]$-module qui est
de type fini et projectif en tant que $\Lambda$-module.

\medskip\noindent
Notre but est de calculer la trace de $1 \in \mathbf{N}$ agissant sur le $\Lambda$-module
des coinvariants de $G$ dans $P \otimes_A M$, c'est-à-dire l'élément
$$
\text{Tr}_{\Lambda}\left(1; \left(P \otimes_A M\right)_G\right) \in
\Lambda^\natural.
$$
L'élément $1\in \mathbf{N}$ correspondra au Frobenius.

\begin{lemma}
\label{lemma-epsilon}
Soit $e\in G$ l'élément neutre. L'application
$$
\begin{matrix}
\Lambda[G] & \longrightarrow & \Lambda^{\natural}\\
\sum \lambda_g\cdot g & \longmapsto & \lambda_e
\end{matrix}
$$
se factorise par $\Lambda[G]^\natural$. On note
$\varepsilon : \Lambda[G]^\natural\to \Lambda^\natural$ l'application induite.
\end{lemma}

\begin{proof}
Il faut montrer que l'application annule les commutateurs. On a
$$
\left(\sum\lambda_g g\right)\left(\sum\mu_g g\right)-\left(\sum \mu_g
g\right)\left(\sum\lambda_g g\right)
= \sum_g\left(\sum_{g_1g_2=g}
\lambda_{g_1}\mu_{g_2}-\mu_{g_1}\lambda_{g_2}\right)g
$$
Le coefficient de $e$ est
$$
\sum_g\left(\lambda_g\mu_{g^{-1}}-\mu_g\lambda_{g^{-1}}\right) =
\sum_g\left(\lambda_g\mu_{g^{-1}}-\mu_{g^{-1}}\lambda_g\right)
$$
qui est une somme de commutateurs et est donc nul dans $\Lambda^\natural$.
\end{proof}

\begin{definition}
\label{definition-trace-G}
Soit $f : P\to P$ un
endomorphisme d'un $\Lambda[G]$-module projectif de type fini
$P$. On pose
$$
\text{Tr}_{\Lambda}^G(f; P) := \varepsilon\left(\text{Tr}_{\Lambda[G]}(f;
P)\right)
$$
et l'on appelle cet élément la {\it $G$-trace de $f$ sur $P$}.
\end{definition}

\begin{lemma}
\label{lemma-lambda-trace}
Soit $f : P\to P$ un endomorphisme du
$\Lambda[G]$-module projectif de type fini $P$. Alors
$$
\text{Tr}_{\Lambda}(f; P) = \# G \cdot \text{Tr}_\Lambda^G(f; P).
$$
\end{lemma}

\begin{proof}
Par additivité, il suffit de traiter le cas $P = \Lambda[G]$.
Dans ce cas, $f$ est donné par
la multiplication à droite par un élément $\sum\lambda_g\cdot g$ de $\Lambda[G]$. Dans
la base $(g)_{g \in G}$, la matrice de $f$ a pour coefficient
$\lambda_{g_2^{-1}g_1}$ à la position $(g_1, g_2)$. En particulier, tous les
coefficients diagonaux sont égaux à $\lambda_e$, et il y en a $\# G$.
\end{proof}

\begin{lemma}
\label{lemma-A-module-structure}
Le morphisme $A\to \Lambda$ définit une structure de $A$-module sur $\Lambda^\natural$.
\end{lemma}

\begin{proof}
C'est clair.
\end{proof}

\begin{lemma}
\label{lemma-diagonal-action-projective-module}
Soient $P$ un $A[G]$-module projectif de type fini et $M$ un $\Lambda[G]$-module,
projectif de type fini en tant que $\Lambda$-module. Alors $P \otimes_A M$ est un
$\Lambda[G]$-module projectif de type fini, pour la structure induite par l'action
diagonale de $G$.
\end{lemma}

\noindent
Notons que $P \otimes_A M$ est naturellement un $\Lambda$-module, puisque $M$ en est un.
Explicitement, avec l'action diagonale, cela s'écrit
$$
\left(\sum\lambda_g g\right)\left(p \otimes m\right)
=
\sum g p \otimes \lambda_g g m.
$$

\begin{proof}
Pour tout $\Lambda[G]$-module $N$, on a
$$
\Hom_{\Lambda[G]}\left(P \otimes_A M, N\right)= \Hom_{A[G]}\left(P,
\Hom_{\Lambda}(M, N)\right)
$$
où l'action de $G$ sur $\Hom_{\Lambda}(M, N)$ est donnée par $(g\cdot
\varphi)(m) = g \varphi (g^{-1} m) $. Il suffit alors de remarquer que le
membre de droite est un composé de foncteurs exacts, grâce à la projectivité
de $P$ et de $M$.
\end{proof}

\begin{lemma}
\label{lemma-multiplicative-trace}
Sous les hypothèses du
lemme \ref{lemma-diagonal-action-projective-module},
soient
$u\in \text{End}_{A[G]}(P)$ et $v\in \text{End}_{\Lambda[G]}(M)$. Alors
$$
\text{Tr}_\Lambda^G \left(u \otimes v; P \otimes_A M\right) = \text{Tr}_A^G(u;
P)\cdot \text{Tr}_\Lambda(v;M).
$$
\end{lemma}

\begin{proof}[Esquisse de démonstration]
On se ramène au cas $P=A[G]$. Dans ce cas, $u$ est la multiplication à droite par un
élément $a = \sum a_gg$ de $A[G]$, et l'on écrit $u = R_a$. Il existe un
isomorphisme de $\Lambda[G]$-modules
$$
\begin{matrix}
\varphi : & A[G]\otimes_A M & \cong & \left(A[G]\otimes_A M\right)'\\
& g \otimes m & \longmapsto & g \otimes g^{-1}m
\end{matrix}
$$
où $\left(A[G]\otimes_A M\right)'$ est muni de la structure de module définie par
l'action à gauche de $G$, ainsi que par la $\Lambda$-linéarité sur $M$. Ce transport
de structure transforme $u \otimes v$ en $\sum_ga_gR_g \otimes g^{-1}v$. En d'autres
termes,
$$
\varphi \circ (u \otimes v) \circ \varphi^{-1}
=
\sum_ga_gR_g \otimes g^{-1}v.
$$
En explicitant les deux membres de l'égalité, il suffit de montrer que
$$
\text{Tr}_\Lambda^G\left(\sum_g a_gR_g \otimes g^{-1}v\right) = a_e\cdot
\text{Tr}_\Lambda(v; M).
$$
Pour cela, on montre que
$$
\text{Tr}_\Lambda^G\left(a_gR_g \otimes g^{-1}v\right) =
\left\{
\begin{matrix}
0 & \text{ si } g\neq e\\
a_e\text{Tr}_\Lambda\left(v; M\right) & \text{ si }g = e
\end{matrix}
\right.
$$
en se ramenant au cas $M=\Lambda$.
\end{proof}

\noindent
Notation :
Considérons l'extension de monoïdes $1 \to G\to \Gamma\to \mathbf{N} \to 1$ et soit
$\gamma\in \Gamma$.
On pose alors $Z_\gamma = \{g\in G | g\gamma = \gamma g\}$.

\begin{lemma}
\label{lemma-gamma-z-gamma-trace}
Soit $P$ un $\Lambda[\Gamma]$-module, projectif de type fini en tant que
$\Lambda[G]$-module, et soit $\gamma \in \Gamma$. Alors
$$
\text{Tr}_{\Lambda}(\gamma, P) =
\# Z_\gamma \cdot \text{Tr}_\Lambda^{Z_\gamma}\left(\gamma, P\right).
$$
\end{lemma}

\begin{proof}
Cela résulte immédiatement du lemme \ref{lemma-lambda-trace}.
\end{proof}

\begin{lemma}
\label{lemma-weak-trace}
Soit $P$ un $A[\Gamma]$-module, projectif de type fini en tant que $A[G]$-module. Soit $M$
un $\Lambda[\Gamma]$-module, projectif de type fini en tant que $\Lambda$-module. Alors
$$
\text{Tr}_{\Lambda}^{Z_\gamma}(\gamma, P \otimes_A M) =
\text{Tr}_A^{Z_\gamma}(\gamma, P)\cdot \text{Tr}_\Lambda(\gamma, M).
$$
\end{lemma}

\begin{proof}
Cela résulte directement du lemme \ref{lemma-multiplicative-trace}.
\end{proof}

\begin{lemma}
\label{lemma-trivial-trace}
Soit $P$ un $\Lambda[\Gamma]$-module, projectif de type fini en tant que
$\Lambda[G]$-module. Alors les coinvariants
$P_G = \Lambda \otimes_{\Lambda[G]} P$
forment un $\Lambda$-module projectif de type fini, muni d'une action de
$\Gamma/G = \mathbf{N}$. De plus, on a
$$
\text{Tr}_\Lambda(1; P_G) =
\sum\nolimits'_{\gamma \mapsto 1} \text{Tr}_\Lambda^{Z_\gamma}(\gamma, P)
$$
où $\sum_{\gamma\mapsto 1}'$ désigne la somme étendue aux classes de $G$-conjugaison
dans $\Gamma$.
\end{lemma}

\begin{proof}[Esquisse de démonstration]
Nous le démontrons d'abord après multiplication par $\# G$.
$$
\# G\cdot \text{Tr}_\Lambda(1; P_G)
= \text{Tr}_\Lambda(\sum\nolimits_{\gamma\mapsto 1} \gamma, P_G)
= \text{Tr}_\Lambda(\sum\nolimits_{\gamma\mapsto 1} \gamma, P)
$$
où la seconde égalité résulte du triangle commutatif
$$
\xymatrix{
P^G \ar[rd]_a & & P_G \ar[ll]^c \\
& P \ar[ur]_b
}
$$
où $a$ est l'inclusion canonique, $b$ la surjection canonique et $c =
\sum_{\gamma \mapsto 1} \gamma$. On a alors
$$
(\sum\nolimits_{\gamma \mapsto 1} \gamma) |_P = a \circ c \circ b
\quad\text{et}\quad
(\sum\nolimits_{\gamma \mapsto 1} \gamma) |_{P_G} = b \circ a \circ c
$$
ils ont donc même trace. On a alors
$$
\# G\cdot \text{Tr}_\Lambda(1; P_G)
=
{\sum_{\gamma\mapsto 1}}'
\frac{\# G}{\# Z_\gamma}\text{Tr}_\Lambda(\gamma, P)
= \# G{\sum_{\gamma\mapsto 1}}' \text{Tr}_\Lambda^{Z_\gamma}(\gamma, P).
$$
Pour achever la démonstration, on se ramène au cas où $\Lambda$ est sans torsion par un
argument d'universalité. Voir \cite{SGA4.5} pour les détails.
\end{proof}

\begin{remark}
\label{remark-content-trivial-trace}
Illustrons le contenu de la formule du
lemme \ref{lemma-weak-trace}.
Supposons que $\Lambda$, vu comme $\Gamma$-module trivial, admette une
résolution finie
$
0\to P_r\to \ldots \to P_1 \to P_0\to \Lambda\to 0
$
par des $\Lambda[\Gamma]$-modules $P_i$ qui sont projectifs de type fini en tant que
$\Lambda[G]$-modules. Dans ce cas,
$$
H_*\left(\left(P_\bullet\right)_G\right) =
\text{Tor}_*^{\Lambda[G]}\left(\Lambda, \Lambda\right) = H_*(G, \Lambda)
$$
et
$$
\text{Tr}_\Lambda^{Z_\gamma}\left(\gamma, P_\bullet\right) =\frac{1}{\#
Z_\gamma}\text{Tr}_\Lambda(\gamma, P_\bullet)=\frac{1}{\#
Z_\gamma}\text{Tr}(\gamma, \Lambda) = \frac{1}{\# Z_\gamma}.
$$
Par conséquent, le lemme \ref{lemma-weak-trace} donne
$$
\text{Tr}_\Lambda (1 , P_G)
= \text{Tr}\left(1 |_{H_*(G, \Lambda)}\right)
= {\sum_{\gamma\mapsto 1}}'\frac{1}{\# Z_\gamma}.
$$
On peut interpréter ceci comme un comptage de points sur le champ $BG$. Si
$\Lambda = \mathbf{F}_\ell$ avec $\ell$ premier à $\# G$, alors
$H_*(G, \Lambda)$ est égal à $\mathbf{F}_\ell$ en degré 0 (et à $0$ aux
autres degrés) et la formule s'écrit
$$
1 =
\sum\nolimits_{
\frac{\sigma\text{-conjugaison}}{\text{classes}\langle\gamma\rangle}
}
\frac{1}{\# Z_\gamma} \mod \ell.
$$
C'est, en un certain sens, une formule des traces ``triviale'' pour $G$.
Nous verrons plus loin que (\ref{equation-trace-formula}) peut, dans
certains cas, être considérée comme une formule des traces hautement non triviale pour un
certain type de groupe ; voir la
section \ref{section-abstract-trace-formula}.
\end{remark}





%11.12.09
\section{Démonstration de la formule des traces}
\label{section-proof-trace-formula}

\begin{theorem}
\label{theorem-trace-formula-again}
Soit $k$ un corps fini et soit $X$ un schéma séparé de type fini de dimension
au plus 1 sur $k$. Soit $\Lambda$ un anneau fini dont le cardinal est premier
à celui de $k$, et soit $K\in D_{ctf}(X, \Lambda)$. Alors
\begin{equation}
\label{equation-trace-formula-again}
\text{Tr}(\pi_X^* |_{R\Gamma_c(X_{\bar k}, K)})
=
\sum\nolimits_{x\in X(k)}
\text{Tr}(\pi_x |_{K_{\bar x}})
\end{equation}
dans $\Lambda^{\natural}$.
\end{theorem}

\noindent
On trouvera dans la remarque \ref{remark-on-trace-formula-again} quelques remarques
sur l'énoncé. Notation : pour abréger, nous écrivons
$$
T'(X, K) =
\sum\nolimits_{x\in X(k)}
\text{Tr}(\pi_x |_{K_{\bar x}})
$$
pour le membre de droite de (\ref{equation-trace-formula-again}) et
$$
T''(X, K)
=\text{Tr}(\pi_x^* |_{R\Gamma_c(X_{\bar k}, K)})
$$
pour le membre de gauche.

\begin{proof}[Démonstration du théorème \ref{theorem-trace-formula-again}]
La démonstration se fait en plusieurs étapes.

\medskip\noindent
Étape 1. {\it Soit $j : \mathcal{U}\hookrightarrow X$ une immersion ouverte de
complément $Y = X - \mathcal{U}$ et soit $i : Y \hookrightarrow X$. Alors
$T''(X, K) = T''(\mathcal{U}, j^{-1} K)+ T''(Y, i^{-1}K)$ et
$T'(X, K) = T'(\mathcal{U}, j^{-1} K)+ T'(Y, i^{-1}K)$.}

\medskip\noindent
C'est clair pour $T'$. Pour $T''$, utilisons la suite exacte
$$
0\to j_!j^{-1} K \to K \to i_* i^{-1} K \to 0
$$
pour munir $K$ d'une filtration. On obtient ainsi un objet
$\widetilde K \in DF(X, \Lambda)$
dont les gradués sont $j_!j^{-1}K$ et $i_*i^{-1}K$;
tous deux appartiennent à $D_{ctf}(X, \Lambda)$. Par les généralités abstraites sur les catégories dérivées filtrées
(INSÉRER UNE RÉFÉRENCE),
$R\Gamma_c(X_{\bar k}, K)\in DF_{perf}(\Lambda)$,
et cet objet est muni de $\pi_x^*$ dans
$DF_{perf}(\Lambda)$.
D'après l'analyse des traces sur les complexes filtrés (INSÉRER UNE RÉFÉRENCE), on obtient
\begin{eqnarray*}
\text{Tr}(\pi_X^* |_{R\Gamma_c(X_{\bar k}, K)})
& = & \text{Tr}(\pi_X^* |_{R\Gamma_c(X_{\bar k}, j_!j^{-1}K)}) +
\text{Tr}(\pi_X^* |_{R\Gamma_c(X_{\bar k}, i_*i^{-1}K)}) \\
& = & T''(U, i^{-1}K) + T''(Y, i^{-1}K).
\end{eqnarray*}

\noindent
Étape 2. {\it Le théorème est vrai si $\dim X\leq 0$. }

\medskip\noindent
En effet, dans ce cas
$$
R\Gamma_c(X_{\bar k}, K) = R\Gamma(X_{\bar k}, K) = \Gamma(X_{\bar k}, K) =
\bigoplus\nolimits_{\bar x \in X_{\bar k}} K_{\bar x}
$$
qui est muni de l'endomorphisme $\pi_X^*$.
Comme les points fixes de $\pi_X : X_{\bar k}\to X_{\bar k}$ sont exactement les
points $\bar x \in X_{\bar k}$ situés au-dessus d'un point $k$-rationnel $x \in X(k)$
nous obtenons
$$
\text{Tr}\big(\pi_X^*|_{R\Gamma_c(X_{\bar k}, K)}\big) =
\sum\nolimits_{x \in X(k)} \text{Tr}(\pi_x|_{K_{\bar x}}).
$$
ce qui est le résultat voulu.

\medskip\noindent
Étape 3. {\it Il suffit de démontrer l'égalité
$T'(\mathcal{U}, \mathcal{F}) = T''(\mathcal{U}, \mathcal{F})$
dans le cas où
\begin{itemize}
\item $\mathcal{U}$ est une courbe affine irréductible lisse sur $k$,
\item $\mathcal{U}(k) = \emptyset$,
\item $K=\mathcal{F}$ est un faisceau fini localement constant de $\Lambda$-modules
sur $\mathcal{U}$ dont les fibres sont des $\Lambda$-modules projectifs de type fini, et
\item $\Lambda$ est annulé par une puissance d'un nombre premier $\ell$ et $\ell \in k^*$.
\end{itemize}
}

\medskip\noindent
En effet, d'après l'étape 2, nous pouvons retrancher tout ensemble fini de points. Mais nous
n'avons qu'un nombre fini de points rationnels, de sorte que nous pouvons supposer qu'il n'y en a
aucun\footnote{À ce stade, un rire maléfique devrait retentir à l'arrière-plan.}.
Nous pouvons supposer que $\mathcal{U}$ est lisse, irréductible et affine en passant aux
composantes irréductibles et en retranchant au besoin les mauvais points. Les
hypothèses sur $\mathcal{F}$ s'obtiennent en explicitant la définition de
$D_{ctf}(X, \Lambda)$ et celles sur $\Lambda$ de l'examen de sa décomposition
primaire.

\medskip\noindent
Dans la suite de la démonstration, nous considérons la situation
$$
\xymatrix{
\mathcal{V} \ar[d]_f \ar[r] & Y \ar[d]^{\bar f} \\
\mathcal{U} \ar[r] & X
}
$$
où $\mathcal{U}$ est comme ci-dessus, $f$ est un revêtement \'etale galoisien fini,
$\mathcal{V}$ est connexe et les flèches horizontales sont des
compactifications projectives. En posant $G=\text{Aut}(\mathcal{V}|\mathcal{U})$, nous supposons également
(comme nous le pouvons) que $f^{-1}\mathcal{F} =\underline M$ est constant, où le
module $M = \Gamma(\mathcal{V}, f^{-1}\mathcal{F})$ est un $\Lambda[G]$-module
projectif de type fini sur $\Lambda$. Cela correspond à l'extension de
monoïdes triviale
$$
1\to G\to \Gamma = G \times \mathbf{N}\to \mathbf{N}\to 1.
$$
Dans ce contexte, compte tenu des réductions ci-dessus, nous devons montrer que
$T''(\mathcal{U}, \mathcal{F}) = 0$.

\medskip\noindent
Étape 4. {\it Il existe une action naturelle de $G$ sur $f_*f^{-1}\mathcal{F}$ et
le morphisme trace $f_*f^{-1}\mathcal{F}\to \mathcal{F}$ définit un isomorphisme}
$$
(f_*f^{-1}\mathcal{F})\otimes_{\Lambda[G]} \Lambda =
(f_*f^{-1}\mathcal{F})_G \cong \mathcal{F}.
$$

\medskip\noindent
Pour le démontrer, il suffit de tout expliciter en un point géométrique.

\medskip\noindent
Étape 5. {\it Soit $A = \mathbf{Z}/\ell^n \mathbf{Z}$ avec $n\gg 0$. Alors
$f_*f^{-1}\mathcal{F} \cong (f_*\underline A)
\otimes_{\underline A} \underline M$ muni de l'action diagonale de $G$.}

\medskip\noindent
Étape 6. {\it Il existe un isomorphisme canonique
$(f_*\underline A \otimes_{\underline A} \underline M)
\otimes_{\Lambda[G]} \underline \Lambda \cong \mathcal{F}$.
}

\medskip\noindent
En fait, il s'agit d'un produit tensoriel dérivé, grâce à l'hypothèse de projectivité
faite sur $\mathcal{F}$.

\medskip\noindent
Étape 7. {\it Il existe un isomorphisme canonique
$$
R\Gamma_c(\mathcal{U}_{\bar k}, \mathcal{F})
= (R\Gamma_c(\mathcal{U}_{\bar k}, f_*A)\otimes_A^\mathbf{L}
M)\otimes_{\Lambda[G]}^\mathbf{L} \Lambda,
$$
compatible avec l'action de $\pi^*_\mathcal{U}$.
}

\medskip\noindent
Cela résulte du théorème des coefficients universels, c'est-à-dire du fait que
$R\Gamma_c$ commute à $\otimes^\mathbf{L}$, et de la platitude de
$\mathcal{F}$ comme $\Lambda$-module.

\medskip\noindent
Nous avons
\begin{eqnarray*}
\text{Tr}(
\pi_\mathcal{U}^* |_{R\Gamma_c(\mathcal{U}_{\bar k}, \mathcal{F})})
& = &
{\sum_{g \in G}}'
\text{Tr}_{\Lambda}^{Z_g}
\left(
(g, \pi_\mathcal{U}^*)
|_{R\Gamma_c(\mathcal{U}_{\bar k}, f_*A)\otimes_A^\mathbf{L} M}
\right) \\
& = &
{\sum_{g\in G}}'
\text{Tr}_A^{Z_g}
(
(g, \pi_\mathcal{U}^*) |_{R\Gamma_c(\mathcal{U}_{\bar k}, f_*A)}
)
\cdot
\text{Tr}_\Lambda(g|_M)
\end{eqnarray*}
où l'action de $\Gamma$ sur $R\Gamma_c(\mathcal{U}_{\bar k}, \mathcal{F})$ prolonge celle de $G$
et $(e, 1)$ agit par $\pi_\mathcal{U}^*$. L'extension de monoïdes est donc donnée
par $\Gamma = G \times \mathbf{N} \to \mathbf{N}$, $\gamma \mapsto 1$. La première
égalité résulte du lemme \ref{lemma-trivial-trace} et la seconde du
lemme \ref{lemma-weak-trace}.

\medskip\noindent
Étape 8. {\it Il suffit de montrer que
$\text{Tr}_A^{Z_g}((g, \pi_\mathcal{U}^*)
|_{R\Gamma_c(\mathcal{U}_{\bar k}, f_*A)}) \in A$
s'envoie sur zéro dans $\Lambda$.
}

\medskip\noindent
Rappelons que
\begin{eqnarray*}
\# Z_g \cdot \text{Tr}_A^{Z_g}((g, \pi_\mathcal{U}^*)
|_{R\Gamma_c(\mathcal{U}_{\bar k}, f_*A)})
& = & \text{Tr}_A((g, \pi_\mathcal{U}^*)
|_{R\Gamma_c(\mathcal{U}_{\bar k}, f_*A)})\\
& = &
\text{Tr}_A((g^{-1}\pi_\mathcal{V})^* |_{R\Gamma_c(\mathcal{V}_{\bar k}, A)}).
\end{eqnarray*}
La première égalité découle du
le lemme \ref{lemma-gamma-z-gamma-trace},
la seconde de la suite spectrale de Leray,
en utilisant la finitude de $f$ et le fait que nous ne prenons que
des traces sur $A$. Or, comme $A=\mathbf{Z}/\ell^n\mathbf{Z}$ avec
$n \gg 0$ et $\# Z_g = \ell^a$ pour un certain entier (fixé) $a$,
il suffit de montrer le résultat suivant.

\medskip\noindent
Étape 9. {\it On a
$\text{Tr}_A((g^{-1}\pi_\mathcal{V})^* |_{R\Gamma_c(\mathcal{V}, A)}) = 0$
dans $A$.}

\medskip\noindent
À nouveau, par additivité, nous avons
\begin{eqnarray*}
& \text{Tr}_A((g^{-1}\pi_\mathcal{V})^* |_{R\Gamma_c(\mathcal{V}_{\bar k} A)})
+
\text{Tr}_A((g^{-1}\pi_\mathcal{V})^*
|_{R\Gamma_c(Y-\mathcal {V})_{\bar k}, A)}) \\
& =
\text{Tr}_A((g^{-1}\pi_Y)^* |_{R\Gamma(Y_{\bar k}, A)})
\end{eqnarray*}
Cette dernière trace est le nombre de points fixes de $g^{-1}\pi_Y$ sur $Y$, d'après
la formule des traces de Weil,
théorème \ref{theorem-weil-trace-formula}.
De plus, d'après le cas de dimension 0 déjà démontré à l'étape 2,
$$
\text{Tr}_A((g^{-1}\pi_\mathcal{V})^*|_{R\Gamma_c(Y-\mathcal{V})_{\bar k}, A)})
$$
est le nombre de points fixes de $g^{-1}\pi_Y$ sur $(Y-\mathcal{V})_{\bar k}$.
Par conséquent,
$$
\text{Tr}_A((g^{-1}\pi_\mathcal{V})^* |_{R\Gamma_c(\mathcal{V}_{\bar k}, A)})
$$
est le nombre de points fixes de $g^{-1}\pi_Y$ sur $\mathcal{V}_{\bar k}$. Mais
il n'y en a aucun : si $\bar y\in Y_{\bar k}$ est fixe par
$g^{-1}\pi_Y$, alors $\bar f(\bar y) \in X_{\bar k}$ est fixe par $\pi_X$. Mais
$\mathcal{U}$ n'a aucun point $k$-rationnel, donc on doit avoir $\bar f(\bar y)\in
(X-\mathcal{U})_{\bar k}$ et ainsi $\bar y\notin \mathcal{V}_{\bar k}$, une
contradiction.
Ceci achève la démonstration.
\end{proof}

\begin{remark}
\label{remark-on-trace-formula-again}
Remarques sur le théorème \ref{theorem-trace-formula-again}.
\begin{enumerate}
\item
Cette formule vaut en toute dimension. Par un lemme de d\'evissage (qui utilise le changement de
base propre, etc.), elle se ramène, même dans cette généralité, à l'énoncé actuel.
\item
Le complexe $R\Gamma_c(X_{\bar k}, K)$ est défini en choisissant une immersion ouverte
$j : X \hookrightarrow \bar X$ avec $\bar X$ projectif sur $k$ et de dimension au
plus 1, et en posant
$$
R\Gamma_c(X_{\bar k}, K) := R\Gamma(\bar X_{\bar k}, j_!K).
$$
L'indépendance de ce choix de $\bar X$ résulte de
(insérer ici une référence). Nous définissons $H^i_c(X_{\bar k}, K)$
comme le $i$-ième groupe de cohomologie de $R\Gamma_c(X_{\bar k}, K)$.
\end{enumerate}
\end{remark}

\begin{remark}
\label{remark-stronger}
Bien que nous n'ayons fait que des réductions et surtout de l'algèbre, la formule des traces
du théorème \ref{theorem-trace-formula-again} est bien plus forte que
la formule géométrique des traces de Weil (théorème \ref{theorem-weil-trace-formula})
car elle s'applique aux systèmes de coefficients
(aux faisceaux), et pas seulement aux coefficients constants.
\end{remark}


%11.17.09
\section{Applications}
\label{section-applications}

\noindent
Bon, après avoir indiqué la démonstration de la formule des traces, essayons d'en tirer
quelque chose.





\section{Sur les faisceaux l-adiques}
\label{section-l-adic-sheaves}

\begin{definition}
\label{definition-l-adic-sheaf}
Soit $X$ un schéma noethérien. Un {\it faisceau de $\mathbf{Z}_\ell$-modules} sur $X$, ou
simplement un {\it faisceau $\ell$-adique} $\mathcal{F}$, est un
système projectif $\left\{\mathcal{F}_n\right\}_{n\geq 1}$ dans lequel
\begin{enumerate}
\item
$\mathcal{F}_n$ est un faisceau constructible de $\mathbf{Z}/\ell^n\mathbf{Z}$-modules sur
$X_\etale$, et
\item
les morphismes de transition $\mathcal{F}_{n+1}\to \mathcal{F}_n$ induisent des isomorphismes
$\mathcal{F}_{n+1} \otimes_{\mathbf{Z}/\ell^{n+1}\mathbf{Z}}
\mathbf{Z}/\ell^n\mathbf{Z} \cong \mathcal{F}_n$.
\end{enumerate}
On dit que $\mathcal{F}$ est {\it lisse} si chaque $\mathcal{F}_n$ est localement
constant. Un {\it morphisme} entre de tels faisceaux n'est autre qu'un morphisme de systèmes projectifs.
\end{definition}

\begin{lemma}
\label{lemma-eventually-constant}
Soit $\{\mathcal{G}_n\}_{n\geq 1}$ un système projectif de faisceaux constructibles de
$\mathbf{Z}/\ell^n\mathbf{Z}$-modules.
Supposons que, pour tout $k\geq 1$, les morphismes
$$
\mathcal{G}_{n+1}/\ell^k \mathcal{G}_{n+1}\to \mathcal{G}_n /\ell^k
\mathcal{G}_n
$$
soient des isomorphismes pour tout $n\gg 0$ (la borne pouvant dépendre de $k$).
En d'autres termes, supposons que le système
$\{\mathcal{G}_n/\ell^k\mathcal{G}_n\}_{n\geq 1}$
soit essentiellement constant, et désignons par $\mathcal{F}_k$ le faisceau correspondant.
Alors le système $\left\{\mathcal{F}_k\right\}_{k\geq 1}$ est un
faisceau de $\mathbf{Z}_\ell$-modules sur $X$.
\end{lemma}

\begin{proof}
La démonstration est immédiate.
\end{proof}

\begin{lemma}
\label{lemma-l-adic-abelian}
La catégorie des faisceaux de $\mathbf{Z}_\ell$-modules sur $X$ est abélienne.
\end{lemma}

\begin{proof}
Soit
$\Phi = \left\{\varphi_n\right\}_{n\geq 1} :
\left\{\mathcal{F}_n\right\}
\to
\left\{\mathcal{G}_n\right\}$
un morphisme de faisceaux de $\mathbf{Z}_\ell$-modules. Posons
$$
\Coker(\Phi) =
\left\{
\Coker\left(\mathcal{F}_n \xrightarrow{\varphi_n} \mathcal{G}_n\right)
\right\}_{n\geq 1}
$$
et $\Ker(\Phi)$ s'obtient en appliquant le
lemme \ref{lemma-eventually-constant}
au système projectif
$$
\left\{
\bigcap_{m\geq n}
\Im\left(\Ker(\varphi_m) \to \Ker(\varphi_n)\right)
\right\}_{n \geq 1}.
$$
On laisse au lecteur le soin de vérifier que ceci définit une catégorie abélienne.
\end{proof}

\begin{example}
\label{example-kernel}
Soit $X=\Spec(\mathbf{C})$ et soit $\Phi : \mathbf{Z}_\ell\to \mathbf{Z}_\ell$
la multiplication par $\ell$. Plus précisément,
$$
\Phi = \left\{ \mathbf{Z}/\ell^n\mathbf{Z} \xrightarrow{\ell}
\mathbf{Z}/\ell^n\mathbf{Z}\right\}_{n \geq 1}.
$$
Pour calculer le noyau, considérons le système projectif
$$
\ldots\to \mathbf{Z}/\ell\mathbf{Z}\xrightarrow{0}
\mathbf{Z}/\ell\mathbf{Z}\xrightarrow{0}\mathbf{Z}/\ell\mathbf{Z}.
$$
Puisque les images sont toujours nulles, $\Ker(\Phi)$ est le système nul.
\end{example}

\begin{remark}
\label{remark-stalk-l-adic-sheaf}
Si $\mathcal{F} = \left\{\mathcal{F}_n\right\}_{n\geq 1}$ est un
faisceau de $\mathbf{Z}_\ell$-modules sur $X$ et si $\bar x$ est un point géométrique, alors
$M_n = \left\{\mathcal{F}_{n, \bar x}\right\}$ est un système projectif de
$\mathbf{Z}/\ell^n\mathbf{Z}$-modules de type fini dont $M_{n+1}\to M_n$ est surjectif
et $M_n = M_{n+1}/\ell^n M_{n+1}$. Il en résulte que
$$
M = \lim_n M_n = \lim \mathcal{F}_{n, \bar x}
$$
est un module de type fini sur $\mathbf{Z}_\ell$. Cela résulte des
lemmes \ref{algebra-lemma-limit-complete} et
\ref{algebra-lemma-finite-over-complete-ring} du chapitre Algèbre, ainsi que du fait que
$M/\ell M = M_1$ est de dimension finie sur $\mathbf{F}_\ell$.
On a donc $M\cong \mathbf{Z}_\ell^{\oplus r} \oplus \oplus_{i = 1}^t
\mathbf{Z}_\ell/\ell^{e_i}\mathbf{Z}_\ell$ pour certains $r, t\geq 0$, $e_i\geq 1$.
Le module $M = \mathcal{F}_{\bar x}$ est appelé la {\it fibre} de
$\mathcal{F}$ en $\bar x$.
\end{remark}

\begin{definition}
\label{definition-torsion-l-adic-sheaf}
Un faisceau de $\mathbf{Z}_\ell$-modules $\mathcal{F}$ est {\it de torsion} si
$\ell^n : \mathcal{F} \to \mathcal{F}$ est le morphisme nul pour un certain $n$.
La catégorie abélienne
des faisceaux de $\mathbf{Q}_\ell$-modules sur $X$ est le quotient de la catégorie abélienne des
faisceaux de $\mathbf{Z}_\ell$-modules par la sous-catégorie de Serre des faisceaux de torsion. En
d'autres termes, ses objets sont les faisceaux de $\mathbf{Z}_\ell$-modules sur $X$ et, si
$\mathcal{F}, \mathcal{G}$ sont deux tels faisceaux, alors
$$
\Hom_{\mathbf{Q}_\ell} \left(\mathcal{F}, \mathcal{G} \right) =
\Hom_{\mathbf{Z}_\ell} \left(\mathcal{F}, \mathcal{G}\right)
\otimes_{\mathbf{Z}_\ell} \mathbf{Q}_\ell.
$$
On désigne par $\mathcal{F} \mapsto \mathcal{F} \otimes \mathbf{Q}_\ell$ le
foncteur de passage au quotient (adjoint à droite du foncteur d'inclusion). Si $\mathcal{F} =
\mathcal{F}' \otimes \mathbf{Q}_\ell$ où $\mathcal{F}'$ est un
faisceau de $\mathbf{Z}_\ell$-modules et $\bar x$ un point géométrique, alors la
{\it fibre} de $\mathcal{F}$ en $\bar x$ est $\mathcal{F}_{\bar x} =
\mathcal{F}'_{\bar x} \otimes \mathbf{Q}_\ell$.
\end{definition}

\begin{remark}
\label{remark-torsion-stalks}
Puisqu'un faisceau de $\mathbf{Z}_\ell$-modules n'est défini que sur un schéma noethérien, il est
de torsion si et seulement si ses fibres sont de torsion.
\end{remark}

\begin{definition}
\label{definition-cohomology-l-adic}
Si $X$ est un schéma séparé de type fini sur un corps algébriquement clos
$k$ et si $\mathcal{F} = \left\{\mathcal{F}_n\right\}_{n\geq 1}$ est un
faisceau de $\mathbf{Z}_\ell$-modules sur $X$, on définit
$$
H^i(X, \mathcal{F}) := \lim_n H^i(X, \mathcal{F}_n)
\quad\text{et}\quad
H_c^i(X, \mathcal{F}) := \lim_n H_c^i(X, \mathcal{F}_n).
$$
Si $\mathcal{F} = \mathcal{F}'\otimes \mathbf{Q}_\ell$ pour un
faisceau de $\mathbf{Z}_\ell$-modules $\mathcal{F}'$, on pose
$$
H_c^i(X , \mathcal{F}) := H_c^i(X,
\mathcal{F}')\otimes_{\mathbf{Z}_\ell}\mathbf{Q}_\ell.
$$
On appelle ces groupes la {\it cohomologie $\ell$-adique} de $X$ à coefficients dans
$\mathcal{F}$.
\end{definition}





\section{Fonctions L}
\label{section-L-function}

\begin{definition}
\label{definition-L-function-finite-ring}
Soit $X$ un schéma de type fini sur un corps fini $k$. Soit $\Lambda$ un
anneau fini d'ordre premier à la caractéristique de $k$ et $\mathcal{F}$ un
faisceau constructible et plat de $\Lambda$-modules sur $X_\etale$. On pose alors
$$
L(X, \mathcal{F}) :=
\prod\nolimits_{x\in |X|}
\det(1 - \pi_x^*T^{\deg x} |_{\mathcal{F}_{\bar x}})^{-1} \in \Lambda [[ T ]]
$$
où $|X|$ est l'ensemble des points fermés de $X$, $\deg x = [\kappa(x): k]$ et
$\bar x$ est un point géométrique au-dessus de $x$. Cette définition se
généralise clairement au cas où l'on remplace $\mathcal{F}$ par
$K \in D_{ctf}(X, \Lambda)$. On appelle cette expression la {\it fonction $L$ de
$\mathcal{F}$}.
\end{definition}

\begin{remark}
\label{remark-T}
Intuitivement, il convient de considérer $T$ comme $T = t^f$, où $p^f = \# k$. Les
définitions ne dépendent alors pas du cardinal du corps de base.
\end{remark}

\begin{definition}
\label{definition-L-function-l-adic}
Supposons maintenant que $\mathcal{F}$ soit un faisceau de $\mathbf{Q}_\ell$-modules sur $X$.
Dans ce cas, on définit
$$
L(X, \mathcal{F}) :=
\prod\nolimits_{x \in |X|}
\det(1 - \pi_x^*T^{\deg x} |_{\mathcal{F}_{\bar x}})^{-1}
\in \mathbf{Q}_\ell[[T]].
$$
Notons que ce produit converge puisqu'il n'y a qu'un nombre fini de points de
degré donné. On appelle ce produit la {\it fonction $L$ de
$\mathcal{F}$}.
\end{definition}




\section{Interprétation cohomologique}
\label{section-L-cohomological}

\noindent
C'est ainsi que Grothendieck a interprété la fonction $L$.

\begin{theorem}[Coefficients finis]
\label{theorem-A}
Soit $X$ un schéma de type fini sur un corps fini $k$. Soit $\Lambda$ un
anneau fini d'ordre premier à la caractéristique de $k$ et $\mathcal{F}$ un
faisceau constructible et plat de $\Lambda$-modules sur $X_\etale$. Alors
$$
L(X, \mathcal{F}) =
\det(1 - \pi_X^*\ T |_{R\Gamma_c(X_{\bar k}, \mathcal{F})})^{-1}
\in \Lambda[[T]].
$$
\end{theorem}

\begin{proof}
Démonstration omise.
\end{proof}

\noindent
À ce stade, nous ne savons même pas si chacun des groupes de cohomologie
$H^i_c(X_{\bar k}, \mathcal{F})$ est libre.

\begin{theorem}[Faisceaux adiques]
\label{theorem-B}
Soit $X$ un schéma de type fini sur un corps fini $k$, et soit $\mathcal{F}$ un
faisceau de $\mathbf{Q}_\ell$-modules sur $X$. Alors
$$
L(X, \mathcal{F}) =
\prod\nolimits_i
\det(1 - \pi_X^*T |_{H_c^i(X_{\bar k} , \mathcal{F})})^{(-1)^{i + 1}}
\in \mathbf{Q}_\ell[[T]].
$$
\end{theorem}

\begin{proof}
La démonstration est esquissée ci-dessous.
\end{proof}

\begin{remark}
\label{remark-which-is-harder}
Comme nous n'avons développé qu'une partie de la théorie des traces, et non
celle des déterminants, le théorème \ref{theorem-A}
est plus difficile à démontrer que le
théorème \ref{theorem-B}.
Nous ne démontrerons que ce dernier; pour le premier, voir \cite{SGA4.5}.
Remarquons aussi qu'il n'existe pas de version plus générale de ce théorème à
coefficients dans $\mathbf{Z}_\ell$, puisqu'il n'y a pas de $\ell$-torsion.
\end{remark}

\noindent
La démonstration du théorème \ref{theorem-B} se ramène à une formule des traces. Puisque
$\mathbf{Q}_\ell$ est de caractéristique 0, il suffit de prouver l'égalité après
avoir pris les dérivées logarithmiques. Plus précisément, nous appliquons $T\frac{d}{dT} \log$
aux deux membres. Nous avons d'une part
\begin{align*}
T \frac{d}{dT} \log L(X, \mathcal{F})
& =
T\frac{d}{dT} \log
\prod_{x \in |X|} \det(1 - \pi_x^* T^{\deg x} |_{\mathcal{F}_{\bar x}})^{-1} \\
& =
\sum_{x \in |X|} T \frac{d}{dT} \log
( \det(1 - \pi_x^* T^{\deg x} |_{\mathcal{F}_{\bar x}})^{-1}) \\
& =
\sum_{x \in |X|} \deg x
\sum_{n \geq 1} \text{Tr}((\pi_x^n)^* |_{\mathcal{F}_{\bar x}}) T^{n\deg x}
\end{align*}
où la dernière égalité résulte de la formule
$$
T\frac{d}{dT}\log\left(\det\left(1-fT|_M\right)^{-1}\right) = \sum_{n\geq 1}
\text{Tr}(f^n|_M)T^n
$$
ce qui vaut pour tout anneau commutatif $\Lambda$ et tout endomorphisme $f$ d'un
$\Lambda$-module projectif de type fini $M$. D'autre part, nous avons
\begin{align*}
& T\frac{d}{dT} \log\left(
\prod\nolimits_i
\det(1-\pi_X^*T |_{H_c^i\left(X_{\bar k} , \mathcal{F}\right)})^{(-1)^{i+1}}
\right) \\
& =
\sum\nolimits_i (-1)^i \sum\nolimits_{n \geq 1}
\text{Tr}\left((\pi_X^n)^* |_{H_c^i(X_{\bar k},\mathcal{F})}\right) T^n
\end{align*}
en appliquant encore une fois la même formule. Maintenant, en comparant les puissances de $T$ et en utilisant la
formule d'inversion de Möbius, nous voyons que le théorème \ref{theorem-B} découle de
l'égalité suivante
$$
\sum_{d | n} d \sum_{x \in |X| \atop \deg x = d}
\text{Tr}((\pi_X^{n/d})^* |_{\mathcal{F}_{\bar x}}) =
\sum_i (-1)^i \text{Tr}((\pi^n_X)^* |_{H^i_c(X_{\bar k}, \mathcal{F})}).
$$
Notons $k_n$ l'extension de degré $n$ de $k$,
$X_n = X \times_{\Spec k} \Spec(k_n)$ et
$_n\mathcal{F} = \mathcal{F}|_{X_n}$, cela se ramène
à
$$
\sum_{x \in X_n(k_n)} \text{Tr}(\pi_X^* |_{_n\mathcal{F}_{\bar x}})
=
\sum_i (-1)^i \text{Tr}((\pi^n_X)^* |_{H^i_c({(X_n)}_{\bar k}, _n\mathcal{F})})
$$
ce qui découle du théorème \ref{theorem-C}.

%11.19.09


\begin{theorem}
\label{theorem-D}
Soit $X/k$ comme ci-dessus, soit $\Lambda$ un anneau fini tel que $\#\Lambda \in k^*$
et soit $K\in D_{ctf}(X, \Lambda)$. Alors $R\Gamma_c(X_{\bar k}, K)\in
D_{perf}(\Lambda)$ et
$$
\sum_{x\in X(k)}\text{Tr}\left(\pi_x |_{K_{\bar x}}\right) =
\text{Tr}\left(\pi_X^* |_{R\Gamma_c(X_{\bar k}, K )}\right).
$$
\end{theorem}

\begin{proof}
Remarquons que nous avons déjà démontré ceci (RÉFÉRENCE) lorsque $\dim X \leq 1$. Le
cas général découle aisément de ce cas et du théorème de changement de base
propre.
\end{proof}

\begin{theorem}
\label{theorem-C}
Soit $X$ un schéma séparé de type fini sur un corps fini $k$ et soit
$\mathcal{F}$ un $\mathbf{Q}_\ell$-faisceau sur $X$. Alors
$\dim_{\mathbf{Q}_\ell}H_c^i(X_{\bar k}, \mathcal{F})$ est finie pour tout $i$,
et ne peut être non nulle que pour $0\leq i \leq 2 \dim X$. De plus, nous avons
$$
\sum_{x\in X(k)} \text{Tr}\left(\pi_x |_{\mathcal{F}_{\bar x}}\right) =
\sum_i (-1)^i\text{Tr}\left(\pi_X^* |_{H_c^i(X_{\bar k}, \mathcal{F})}\right).
$$
\end{theorem}

\begin{proof}
Nous expliquons comment déduire ce résultat du théorème \ref{theorem-D}.
Nous utilisons d'abord quelques arguments de cohomologie \'etale pour ramener la démonstration
à un énoncé algébrique que nous démontrons ensuite.

\medskip\noindent
Soit $\mathcal{F}$ comme dans le théorème. Nous pouvons écrire
$\mathcal{F}$ sous la forme
$\mathcal{F}'\otimes \mathbf{Q}_\ell$ où $\mathcal{F}' =
\left\{\mathcal{F}'_n\right\}$ est un $\mathbf{Z}_\ell$-faisceau sans torsion,
c'est-à-dire que $\ell : \mathcal{F}'\to \mathcal{F}'$ a un noyau nul dans la
catégorie des $\mathbf{Z}_\ell$-faisceaux. Alors chaque $\mathcal{F}_n'$ est un
faisceau constructible et plat de $\mathbf{Z}/\ell^n\mathbf{Z}$-modules sur $X_\etale$, donc
$\mathcal{F}'_n \in D_{ctf}(X, \mathbf{Z}/\ell^n\mathbf{Z})$ et
$\mathcal{F}_{n+1}'
\otimes^{\mathbf{L}}_{\mathbf{Z}/\ell^{n+1}\mathbf{Z}}
\mathbf{Z}/\ell^n\mathbf{Z} = \mathcal{F}_n'$.
Remarquons que la dernière égalité est également valable
pour le produit tensoriel usuel (non dérivé), puisque $\mathcal{F}'_n$ est plat
(il s'agit de la même égalité). Par conséquent,
\begin{enumerate}
\item
le complexe $K_n = R\Gamma_c\left(X_{\bar k}, \mathcal{F}_n'\right)$ est parfait,
et il est muni d'un endomorphisme $\pi_n : K_n\to K_n$ dans
$D(\mathbf{Z}/\ell^n\mathbf{Z})$,
\item
on dispose d'identifications
$$
K_{n+1}
\otimes^{\mathbf{L}}_{\mathbf{Z}/\ell^{n+1}\mathbf{Z}}
\mathbf{Z}/\ell^n\mathbf{Z}
=
K_n
$$
dans $D_{perf}(\mathbf{Z}/\ell^n\mathbf{Z})$, compatibles avec les endomorphismes
$\pi_{n+1}$ et $\pi_n$ (voir \cite[Rapport 4.12]{SGA4.5}),
\item
l'égalité $\text{Tr}\left(\pi_X^* |_{K_n}\right) =
\sum_{x\in X(k)} \text{Tr}\left(\pi_x |_{(\mathcal{F}'_n)_{\bar x}}\right)$
est vérifiée, et
\item
pour tout $x\in X(k)$, les éléments
$\text{Tr}(\pi_x |_{\mathcal{F}'_{n, \bar x}}) \in \mathbf{Z}/\ell^n\mathbf{Z}$
définissent un élément de
$\mathbf{Z}_\ell$ qui est égal à
$\text{Tr}(\pi_x |_{\mathcal{F}_{\bar x}}) \in \mathbf{Q}_\ell$.
\end{enumerate}
Il suffit donc de démontrer le lemme d'algèbre suivant.
\end{proof}

\begin{lemma}
\label{lemma-piece-together}
Supposons donnés
$K_n\in D_{perf}(\mathbf{Z}/\ell^n\mathbf{Z})$, $\pi_n : K_n\to K_n$
et des isomorphismes
$\varphi_n :
K_{n+1} \otimes^\mathbf{L}_{\mathbf{Z}/\ell^{n+1}\mathbf{Z}}
\mathbf{Z}/\ell^n\mathbf{Z}
\to K_n$
compatibles avec $\pi_{n+1}$ et $\pi_n$. Alors
\begin{enumerate}
\item
les éléments $t_n = \text{Tr}(\pi_n |_{K_n})\in \mathbf{Z}/\ell^n\mathbf{Z}$
définissent un élément $t_\infty = \{t_n\}$ de $\mathbf{Z}_\ell$,
\item
le $\mathbf{Z}_\ell$-module $H_\infty^i = \lim_n H^i(k_n)$ est de type fini et
est nul sauf pour un nombre fini d'indices $i$, et
\item
les opérateurs $H^i(\pi_n): H^i(K_n)\to H^i(K_n)$ sont compatibles et définissent
$\pi_\infty^i : H_\infty^i\to H_\infty^i$ qui vérifie
$$
\sum (-1)^i \text{Tr}(
\pi_\infty^i |_{H_\infty^i \otimes_{\mathbf{Z}_\ell}\mathbf{Q}_\ell}) =
t_\infty.
$$
\end{enumerate}
\end{lemma}

\begin{proof}
Puisque $\mathbf{Z}/\ell^n\mathbf{Z}$ est un anneau local et que $K_n$ est parfait, chaque
$K_n$ peut être représenté par un complexe fini $K_n^\bullet$ de
$\mathbf{Z}/\ell^n \mathbf{Z}$-modules libres de type fini tel que l'image du morphisme $K_n^p \to K_n^{p+1}$
soit contenue dans $\ell K_n^{p+1}$. Un tel complexe est en fait
unique à isomorphisme près. De plus, $\pi_n$ peut être représenté par un morphisme de
complexes $\pi_n^\bullet : K_n^\bullet\to K_n^\bullet$ (qui est unique à
homotopie près). De même, l'isomorphisme
$\varphi_n : K_{n+1} \otimes_{\mathbf{Z}/\ell^{n+1}\mathbf{Z}}^{\mathbf{L}}
\mathbf{Z}/\ell^n\mathbf{Z}\to K_n$ est représenté par un morphisme de complexes
$$
\varphi_n^\bullet :
K_{n+1}^\bullet
\otimes_{\mathbf{Z}/\ell^{n+1}\mathbf{Z}}
\mathbf{Z}/\ell^n\mathbf{Z} \to K_n^\bullet.
$$
En fait, $\varphi_n^\bullet$ est un isomorphisme de complexes ; nous voyons donc que
\begin{itemize}
\item
il existe $a, b\in \mathbf{Z}$ indépendants de $n$ tels que $K_n^i = 0$
pour tout $i\notin[a, b]$, et
\item
le rang de $K_n^i$ ne dépend pas de $n$.
\end{itemize}
Par conséquent, le module $K_\infty^i = \lim_n \{K_n^i, \varphi_n^i\}$ est un
$\mathbf{Z}_\ell$-module libre de type fini et $K_\infty^\bullet$ est un complexe fini
de $\mathbf{Z}_\ell$-modules libres de type fini. Par récurrence sur le nombre de termes non nuls,
on peut démontrer que $H^i\left(K_\infty^\bullet\right) = \lim_n
H^i\left(K_n^\bullet\right)$ (ceci est faux pour les complexes non bornés). Nous
en concluons que $H_\infty^i = H^i\left(K_\infty^\bullet\right)$ est un
$\mathbf{Z}_\ell$-module de type fini. Cela démontre {\it ii}. Pour démontrer le reste du
lemme, nous devons remédier au défaut éventuel de commutativité des diagrammes
$$
\xymatrix{
{K_{n+1}^\bullet} \ar[d]_{\pi_{n+1}^\bullet} \ar[r]^{\varphi_n^\bullet} &
{K_n^\bullet} \ar[d]^{\pi_n^\bullet} \\
{K_{n+1}^\bullet} \ar[r]_{\varphi_n^\bullet} & {K_n^\bullet.}
}
$$
Cependant, ce diagramme commute bien dans la catégorie dérivée ; il commute donc
à homotopie près. Nous remplaçons par récurrence $\pi_n^\bullet$ pour $n\geq 2$ par
des morphismes de complexes homotopes rendant ces diagrammes commutatifs. Plus précisément, si $h^i :
K_{n+1}^i \to K_n^{i-1}$ est une homotopie, c'est-à-dire
$$
\pi_n^\bullet \circ \varphi_n^\bullet -
\varphi_n^\bullet \circ \pi_{n + 1}^\bullet = dh + hd,
$$
alors nous choisissons $\tilde h^i : K_{n+1}^i\to K_{n+1}^{i-1}$ relevant $h^i$. Cela est
possible parce que $K_{n+1}^i$ est libre et que le morphisme $K_{n+1}^{i-1}\to K_n^{i-1}$ est
surjectif. Nous remplaçons alors $\pi_n^\bullet$ par $\tilde\pi_n^\bullet$ défini par
$$
\tilde\pi_{n+1}^\bullet = \pi_{n+1}^\bullet + d\tilde h+\tilde hd.
$$
Avec ce choix de $\{\pi_n^\bullet\}$, les diagrammes ci-dessus commutent et les
morphismes sont compatibles et définissent un endomorphisme $\pi_\infty^\bullet =
\lim_n\pi_n^\bullet$ de $K_\infty^\bullet$. La partie {\it i} est alors claire :
les éléments $t_n = \sum(-1)^i \text{Tr}\left(\pi_n^i |_{K_n^i}\right)$
définissent un élément $t_\infty$ de $\mathbf{Z}_\ell$. De plus,
\begin{align*}
t_\infty
& =
\sum (-1)^i \text{Tr}_{\mathbf{Z}_\ell}(\pi_\infty^i |_{K_\infty^i}) \\
& =
\sum (-1)^i \text{Tr}_{\mathbf{Q}_\ell}(
\pi_\infty^i |_{K_\infty^i \otimes_{\mathbf{Z}_\ell}\mathbf{Q}_\ell}) \\
& =
\sum (-1)^i \text{Tr}(
\pi_\infty |_{H^i(K_\infty^\bullet \otimes \mathbf{Q}_\ell)})
\end{align*}
où la dernière égalité découle du fait que $\mathbf{Q}_\ell$ est un
corps, de sorte que le complexe $K_\infty^\bullet \otimes \mathbf{Q}_\ell$ est
quasi-isomorphe à sa cohomologie
$H^i(K_\infty^\bullet \otimes \mathbf{Q}_\ell)$. Celle-ci est également égale à
$H^i(K_\infty^\bullet)\otimes_{\mathbf{Z}}\mathbf{Q}_\ell = H_\infty^i \otimes
\mathbf{Q}_\ell$, ce qui achève la démonstration du lemme ainsi que celle du
théorème \ref{theorem-C}.
\end{proof}




\section{Liste des compléments à apporter ci-dessus}
\label{section-list-skipped}

\noindent
Qu'avons-nous omis de démontrer jusqu'ici dans les cours ?
\begin{enumerate}
\item les courbes et leurs jacobiennes,
\item le théorème de changement de base propre,
\item une présentation insuffisante de $R\Gamma_c$,
\item plus généralement, étant donné un morphisme séparé de type fini $f : X \to S$,
avec $S$ quasi-projectif, une présentation de $Rf_!$ sur les faisceaux \'etales.
\item une présentation de $\otimes^{\mathbf{L}}$
\item une explication des raisons pour lesquelles $R\Gamma_c$ commute à $\otimes^{\mathbf{L}}$
\end{enumerate}






%11.24.09
\section{Exemples de fonctions L}
\label{section-examples-L-functions}

\noindent
Nous utilisons le théorème \ref{theorem-B} pour les courbes afin de donner des exemples de fonctions $L$.





\section{Faisceaux constants}
\label{section-L-function-constant-sheaf}

\noindent
Soit $k$ un corps fini, $X$ une courbe lisse géométriquement irréductible sur
$k$ et $\mathcal{F} = \underline{\mathbf{Q}_\ell}$ le faisceau constant. Si
$\bar x$ est un point géométrique de $X$, le module galoisien
$\mathcal{F}_{\bar x} = \mathbf{Q}_\ell$ est trivial, donc
$$
\det(1-\pi_x^*\ T^{\deg x} |_{\mathcal{F}_{\bar x}})^{-1} =
\frac{1}{1-T^{\deg x}}.
$$
En appliquant le théorème \ref{theorem-B}, nous obtenons
\begin{align*}
L(X, \mathcal{F})
& =
\prod_{i = 0}^2
\det(1 - \pi_X^*T |_{H_c^i(X_{\bar k}, \mathbf{Q}_\ell)})^{(-1)^{i+1}} \\
& =
\frac{\det(1 - \pi_X^*T |_{H_c^1(X_{\bar k}, \mathbf{Q}_\ell)})}{
\det(1 - \pi_X^*T |_{H_c^0(X_{\bar k}, \mathbf{Q}_\ell)})
\cdot \det(1 - \pi_X^*T |_{H_c^2(X_{\bar k}, \mathbf{Q}_\ell)})}.
\end{align*}
Pour calculer ce dernier, distinguons deux cas.


\medskip\noindent
{\bf Cas projectif.}
Supposons que $X$ soit projective, de sorte que $H_c^i(X_{\bar k}, \mathbf{Q}_\ell) =
H^i(X_{\bar k}, \mathbf{Q}_\ell)$, et nous avons
$$
H^i(X_{\bar k}, \mathbf{Q}_\ell) =
\left\{
\begin{matrix}
\mathbf{Q}_\ell & \pi_X^* = 1 & \text{si }i = 0, \\
\mathbf{Q}_\ell^{2g} & \pi_X^* = ? & \text{si }i = 1, \\
\mathbf{Q}_\ell & \pi_X^* = q & \text{si }i = 2.
\end{matrix}
\right.
$$
L'identification de l'action de $\pi_X^*$ sur $H^2$ résulte du
Cohomologie \'etale, lemme \ref{etale-cohomology-lemma-pullback-on-h2-curve}
et du fait que le degré
de $\pi_X$ est $q = \#(k)$.
Nous savons peu de choses sur l'action de $\pi_X^*$ sur la cohomologie de degré 1.
Appelons $\alpha_1, \ldots, \alpha_{2g}$ ses valeurs propres dans
$\bar{\mathbf{Q}}_\ell$. En rassemblant ces informations,
le théorème \ref{theorem-B}
donne l'égalité
$$
\prod\nolimits_{x \in |X|} \frac{1}{1 - T^{\deg x}} =
\frac{\det(1- \pi_X^* T|_{H^1(X_{\bar k}, \mathbf{Q}_\ell)})}{(1-T)(1-qT)} =
\frac{(1 - \alpha_1 T) \ldots (1 - \alpha_{2g}T)}{(1-T)(1-qT)}
$$
d'où nous déduisons le résultat suivant.

\begin{lemma}
\label{lemma-count-points-projective}
Soit $X$ une courbe lisse, projective et géométriquement irréductible
sur un corps fini $k$. Alors
\begin{enumerate}
\item la fonction $L$ notée $L(X, \mathbf{Q}_\ell)$ est rationnelle,
\item les valeurs propres $\alpha_1, \ldots, \alpha_{2g}$ de $\pi_X^*$ sur
$H^1(X_{\bar k}, \mathbf{Q}_\ell)$ sont des entiers algébriques
indépendants de $\ell$,
\item le nombre de points de $X$ rationnels sur $k_n$, où $[k_n : k] = n$, est
$$
\# X(k_n) = 1 - \sum\nolimits_{i = 1}^{2g}\alpha_i^n + q^n,
$$
\item pour tout $i$, $|\alpha_i| < q$.
\end{enumerate}
\end{lemma}

\begin{proof}
La partie (3) est le théorème \ref{theorem-C} appliqué à $\mathcal{F} =
\underline{\mathbf{Q}_\ell}$ sur $X \otimes k_n$. Pour la partie (4), utilisons le
résultat suivant.
\end{proof}

\begin{exercise}
\label{exercise-powers}
Soient $\alpha_1, \ldots, \alpha_n \in \mathbf{C}$. Alors, pour tout secteur angulaire
contenant l'axe réel positif, de la forme $C_\varepsilon = \{ z \in
\mathbf{C} \ | \ |\arg z| < \varepsilon \}$ avec $\varepsilon > 0$, il existe
un entier $k \geq 1$ tel que $\alpha_1^k, \ldots, \alpha_n^k \in
C_\varepsilon$.
\end{exercise}

\noindent
Démontrer alors que $|\alpha_i| \leq q$ pour tout $i$. Puis utiliser des
considérations élémentaires sur les nombres complexes pour démontrer (comme dans la preuve du théorème
des nombres premiers) que $|\alpha_i| < q$. En fait, l'hypothèse de Riemann affirme que
l'on a $|\alpha_i| = \sqrt{q}$ pour tout $i$. Nous y reviendrons plus loin.

\medskip\noindent
{\bf Cas affine.}
Supposons maintenant que $X$ soit affine, disons $X= \bar X-\left\{x_1, \ldots, x_n\right\}$
où $j : X \hookrightarrow \bar X$ est une compactification projective non singulière.
Alors $H_c^0(X_{\bar k}, \mathbf{Q}_\ell) = 0$ et $H_c^2(X_{\bar k},
\mathbf{Q}_\ell) = H^2(\bar X_{\bar k}, \mathbf{Q}_\ell)$, de sorte que
le théorème \ref{theorem-B}
s'écrit
$$
L(X, \mathbf{Q}_\ell) = \prod_{x \in |X|}\frac{1}{1 - T^{\deg x}} =
\frac{\det(1-\pi_X^*T |_{H_c^1(X_{\bar k}, \mathbf{Q}_\ell)})}{1 - qT}.
$$
D'autre part, le cas précédent donne
\begin{eqnarray*}
L(X, \mathbf{Q}_\ell) & = & L(\bar X,
\mathbf{Q}_\ell)\prod_{i = 1}^n\left(1-T^{\deg x_i}\right) \\
& = & \frac{\prod_{i = 1}^n(1-T^{\deg
x_i})\prod_{j = 1}^{2g}(1-\alpha_jT)}{(1-T)(1-qT)}.
\end{eqnarray*}
Par conséquent, nous voyons que $\dim H_c^1(X_{\bar k}, \mathbf{Q}_\ell) =
2g+\sum_{i = 1}^n \deg(x_i)-1$, et que les valeurs propres $\alpha_1, \ldots,
\alpha_{2g}$ de $\pi_{\bar X}^*$ agissant sur la cohomologie de degré 1 sont des racines de
l'unité. Plus précisément, chaque $x_i$ fournit un système complet de racines $\deg(x_i)$-ièmes
de l'unité, et l'on omet une occurrence de 1. Pour le voir directement à l'aide des
faisceaux cohérents, considérons la suite exacte courte sur $\bar X$
$$
0\to j_!\mathbf{Q}_\ell\to \mathbf{Q}_\ell\to\bigoplus_{i = 1}^n
\mathbf{Q}_{\ell, x_i}\to 0.
$$
La suite exacte longue de cohomologie s'écrit
$$
0\to \mathbf{Q}_\ell \to \bigoplus_{i = 1}^n \mathbf{Q}_\ell^{\oplus \deg x_i}
\to H_c^1(X_{\bar k}, \mathbf{Q}_\ell) \to H_c^1(\bar X_{\bar k},
\mathbf{Q}_\ell)\to 0
$$
où l'action du Frobenius sur $\bigoplus_{i = 1}^n \mathbf{Q}_\ell^{\oplus
\deg x_i}$ est donnée par la permutation cyclique de chaque terme ; par ailleurs, $H_c^2(X_{\bar k},
\mathbf{Q}_\ell) = H_c^2(\bar X_{\bar k}, \mathbf{Q}_\ell)$.






\section{La famille de Legendre}
\label{section-legendre-family}

\noindent
Soit $k$ un corps fini de caractéristique impaire,
$X = \Spec(k[\lambda, \frac{1}{\lambda(\lambda - 1)}])$, et
considérons la famille de courbes elliptiques
$f : E \to X$ dans $\mathbf{P}^2_X$, dont l'équation affine est $y^2 =
x(x - 1)(x - \lambda)$. Posons $\mathcal{F} = Rf_*^1\mathbf{Q}_\ell =
\left\{R^1f_*\mathbf{Z}/\ell^n\mathbf{Z}\right\}_{n\geq 1} \otimes
\mathbf{Q}_\ell$. Dans cette situation, on a les propriétés suivantes :
\begin{itemize}
\item pour tout $n \geq 1$, le faisceau $R^1f_*(\mathbf{Z}/\ell^n\mathbf{Z})$ est
fini et localement constant -- en fait, il est libre de rang 2 sur
$\mathbf{Z}/\ell^n\mathbf{Z}$,
\item le système $\{R^1f_*\mathbf{Z}/\ell^n\mathbf{Z}\}_{n\geq 1}$ est un faisceau
$\ell$-adique lisse, et
\item pour tout $x\in |X|$,
$\det(1 - \pi_x\ T^{\deg x} |_{\mathcal{F}_{\bar x}}) =
(1 - \alpha_x T^{\deg x})(1 - \beta_x T^{\deg x})$
où $\alpha_x, \beta_x$ sont les valeurs propres du Frobenius géométrique
de $E_x$ agissant sur $H^1(E_{\bar x}, \mathbf{Q}_\ell)$.
\end{itemize}
Notons que $E_x$ n'est défini que sur $\kappa(x)$ et non sur $k$. La démonstration de
ces faits utilise le théorème de changement de base propre et l'acyclicité locale des
morphismes lisses. Pour plus de détails, voir \cite{SGA4.5}. Il en résulte que
$$
L(E/X) := L(X, \mathcal{F}) = \prod_{x\in |X|}
\frac{1}{(1-\alpha_xT^{\deg x})(1-\beta_xT^{\deg x })} .
$$
En appliquant le théorème \ref{theorem-B}, nous obtenons
$$
L(E/X) =
\prod_{i = 0}^2
\det\left(1 - \pi_X^*T |_{H_c^i(X_{\bar k}, \mathcal{F})}\right)^{(-1)^{i+1}},
$$
et nous voyons en particulier qu'il s'agit d'une fonction rationnelle. De plus, il est
assez facile de montrer que $H_c^0(X_{\bar k}, \mathcal{F}) = H_c^2(X_{\bar
k}, \mathcal{F}) = 0$, de sorte que nous avons simplement
$$
L(E/X) = \det(1 - \pi_X^*T |_{H_c^1(X, \mathcal{F})}).
$$
Pour calculer explicitement ce déterminant, considérons la suite spectrale de Leray
du morphisme propre $f : E \to X$ à coefficients dans $\mathbf{Q}_\ell$, à savoir
$$
H_c^i(X_{\bar k}, R^jf_*\mathbf{Q}_\ell) \Rightarrow H_c^{i+j}(E_{\bar
k}, \mathbf{Q}_\ell)
$$
qui dégénère. Nous avons $f_*\mathbf{Q}_\ell = \mathbf{Q}_\ell$ et
$R^1f_*\mathbf{Q}_\ell = \mathcal{F}$. Le faisceau $R^2f_*\mathbf{Q}_\ell =
\mathbf{Q}_\ell(-1)$ est le {\it twist à la Tate} de $\mathbf{Q}_\ell$, c'est-à-dire
qu'il s'agit du faisceau $\mathbf{Q}_\ell$ sur lequel l'action de Galois est donnée par
la multiplication par $\#\kappa(x)$ dans la fibre en $\bar x$. Il s'ensuit que,
pour tout $n\geq 1$,
\begin{align*}
\# E(k_n)
& =
\sum\nolimits_i (-1)^i
\text{Tr}({\pi_E^n}^* |_{H_c^i(E_{\bar k}, \mathbf{Q}_\ell)}) \\
& =
\sum\nolimits_{i, j} (-1)^{i+j}
\text{Tr}({\pi^n_X}^* |_{H_c^i(X_{\bar k}, R^jf_*\mathbf{Q}_\ell)}) \\
& =
(q^n - 2) +
\text{Tr}({\pi_X^n}^* |_{H_c^1(X_{\bar k}, \mathcal{F})}) + q^n(q^n - 2) \\
& =
q^{2n} - q^n - 2 +
\text{Tr}({\pi_X^n}^* |_{H_c^1(X_{\bar k}, \mathcal{F})})
\end{align*}
où la première égalité résulte du
théorème \ref{theorem-C},
la deuxième de
la suite spectrale de Leray et la troisième du calcul explicite des images directes supérieures
de $\mathbf{Q}_\ell$ par $f$. On pourrait aussi écrire
$$
\# E(k_n) = \sum_{x \in X(k_n)} \# E_x(k_n)
$$
et appliquer la formule des traces à chaque courbe. On peut aussi déterminer le nombre de
points rationnels sur $k_n$ par un simple comptage. La section nulle fournit $q^n -2$ points
(on omet les points où $\lambda = 0, 1$), d'où
$$
\# E(k_n) =
q^n - 2 + \# \{y^2 = x(x - 1)(x - \lambda), \lambda\neq 0, 1\}.
$$
Nous avons maintenant
$$
\begin{matrix}
\# \{y^2 = x(x - 1)(x - \lambda),\ \lambda\neq 0, 1\} \\
\\
\quad =
\# \{y^2 = x(x - 1)(x - \lambda)\text{ dans }\mathbf{A}^3\}
- \# \{y^2 = x^2(x - 1)\} - \# \{y^2 = x(x - 1)^2\}\\
\\
\quad = \# \{\lambda = \frac{-y^2}{x(x - 1)} + x,\ x\neq 0, 1\} +
\# \{y^2 = x(x - 1)(x - \lambda), x = 0, 1\} - 2(q^n - \varepsilon_n) \\
\\
\quad = q^n(q^n - 2)+2q^n - 2(q^n - \varepsilon_n)\\
\\
\quad = q^{2n}-2q^n+2\varepsilon_n
\end{matrix}
$$
où $\varepsilon_n = 1$ si $-1$ est un carré dans $k_n$, et 0 sinon,
c'est-à-dire
$$
\varepsilon_n = \frac{1}{2}\left(1+\left(\frac{-1}{k_n}\right)\right) =
\frac{1}{2}\left(1+(-1)^{\frac{q^n - 1}{2}}\right).
$$
Ainsi, $ \# E(k_n) = q^{2n} - q^n - 2+ 2\varepsilon_n$.
En comparant avec la formule précédente, on obtient
$$
\text{Tr}({\pi_X^n}^* |_{H_c^1(X_{\bar k}, \mathcal{F})}) =
2 \varepsilon_n = 1 + (-1)^{\frac{q^n - 1}{2}},
$$
ce qui implique, par un calcul élémentaire dans les nombres complexes, que si $-1$ est un
carré dans $k_n^*$, alors $\dim H_c^1(X_{\bar k}, \mathcal{F}) = 2$ et les
valeurs propres sont $1$ et $1$. Par conséquent, dans ce cas, nous avons
$$
L(E/X) = (1 - T)^2.
$$




\section{Sommes exponentielles}
\label{section-exponential-sums}

\noindent
Un problème classique de théorie des nombres consiste à évaluer des sommes de la forme
$$
S_{a, b}(p) = \sum_{x\in \mathbf{F}_p - \left\{0, 1\right\}} e^{\frac{2\pi
ix^a(x - 1)^b}{p}}.
$$
Dans notre contexte, on peut l'interpréter comme une somme cohomologique de la manière suivante.
Considérons le schéma de base
$S = \Spec(\mathbf{F}_p[x, \frac{1}{x(x - 1)}])$ et la courbe affine
$f : X \to \mathbf{P}^1-\{0, 1, \infty\}$ sur $S$ donnée par l'équation
$y^{p - 1} = x^a(x - 1)^b$. C'est un revêtement galoisien fini \'etale de groupe
$\mathbf{F}_p^*$ et il existe une décomposition
$$
f_*(\bar{\mathbf{Q}}_\ell^*) =
\bigoplus_{\chi : \mathbf{F}_p^*\to \bar{\mathbf{Q}}_\ell^*} \mathcal{F}_\chi
$$
où $\chi$ parcourt les caractères de $\mathbf{F}_p^*$ et
$\mathcal{F}_\chi$ est un faisceau lisse de rang 1 de $\mathbf{Q}_\ell$-modules sur lequel
$\mathbf{F}_p^*$ agit par $\chi$ sur les fibres. On obtient une décomposition correspondante
$$
H_c^1(X_{\bar k}, \mathbf{Q}_\ell) = \bigoplus_\chi H^1(\mathbf{P}_{\bar
k}^1-\{0, 1, \infty\}, \mathcal{F}_\chi)
$$
et l'interprétation cohomologique de la somme exponentielle est donnée par la
formule des traces appliquée à $\mathcal{F}_\chi$ sur $\mathbf{P}^1 - \{0, 1,
\infty\}$ pour un $\chi$ convenable. Elle s'écrit
$$
S_{a, b}(p) =
-\text{Tr}(\pi_X^*
|_{H^1(\mathbf{P}_{\bar k}^1-\{0, 1, \infty\}, \mathcal{F}_\chi)}).
$$
Le yoga général de Weil suggère qu'il devrait y avoir une certaine compensation dans la
somme. En appliquant (approximativement) la formule de Riemann--Hurwitz, nous voyons que
$$
2g_X-2 \approx -2 (p-1) + 3(p-2) \approx p
$$
donc $g_X\approx p/2$, ce qui suggère également que les composantes en $\chi$ sont petites.



%12.01.09
\section{Formule des traces en termes de groupes fondamentaux}
\label{section-trace-formula-fundamental-group}

\noindent
Dans les sections suivantes, nous reformulons entièrement la formule des traces
en termes du groupe fondamental d'une courbe, sauf lorsque la courbe
se trouve être $\mathbf{P}^1$.




\section{Groupes fondamentaux}
\label{section-fundamental-groups}

\noindent
Ce matériau est présenté plus en détail dans le chapitre consacré aux
groupes fondamentaux. Voir
Groupes fondamentaux, section \ref{pione-section-introduction}.
Soit $X$ un schéma connexe et soit $\overline{x}\to X$ un
point géométrique. Considérons le foncteur
$$
\begin{matrix}
F_{\overline{x}}: &
\text{ schémas finis \'etales } \atop \text{ sur } X &
\longrightarrow & \text{ ensembles finis} \\
&
Y/X &
\longmapsto &
F_{\overline{x}}(Y) =
\left\{\text{ points géométriques }\overline y \atop \text{ de } Y
\text{ au-dessus de }\overline{x}\right\} = Y_{\overline{x}}
\end{matrix}
$$
Posons
$$
\pi_1(X, \overline{x})
=
Aut(F_{\overline{x}})
=
\text{ ensemble des automorphismes du foncteur }F_{\overline{x}}
$$
Notons que, pour tout morphisme fini \'etale $Y \to X$, il existe une action
$$
\pi_1(X, \overline{x}) \times F_{\overline{x}}(Y) \to F_{\overline{x}}(Y)
$$

\begin{definition}
\label{definition-open}
Un sous-groupe de la forme
$\text{Stab}(\overline y\in F_{\overline{x}}(Y))\subset \pi_1(X, \overline{x})$
est dit {\it ouvert}.
\end{definition}

\begin{theorem}[Grothendieck]
\label{theorem-fundamental-group}
Soit $X$ un schéma connexe.
\begin{enumerate}
\item Il existe une topologie sur $\pi_1(X, \overline{x})$ telle que les sous-groupes
ouverts forment un système fondamental de voisinages ouverts de $e\in \pi_1(X, \overline
x)$.
\item Muni de la topologie de (1), le groupe
$\pi_1(X, \overline{x})$ est un groupe profini.
\item Le foncteur
$$
\begin{matrix}
\text{ schémas finis } \atop \text{ \'etales sur }X & \to &
\text{ ensembles finis discrets à action continue } \atop \pi_1(X, \overline{x})\text{-ensembles}\\
Y / X& \mapsto & F_{\overline{x}}(Y) \text{ muni de son action naturelle}
\end{matrix}
$$
est une équivalence de catégories.
\end{enumerate}
\end{theorem}

\begin{proof}
Cf. \cite{SGA1}.
\end{proof}

\begin{proposition}
\label{proposition-integral-normal-fundamental-group}
Soit $X$ un schéma noethérien normal intègre. Soit
$\overline y\to X$ un point géométrique algébrique situé
au-dessus du point générique $\eta\in X$. Alors
$$
\pi_x(X, \overline \eta) = Gal(M/\kappa(\eta))
$$
($\kappa(\eta)$ est le corps des fonctions de $X$), où
$$
\kappa(\overline \eta)\supset M\supset \kappa(\eta) = k(X)
$$
est la sous-extension maximale telle que, pour toute sous-extension finie
$M\supset L\supset \kappa(\eta)$, la normalisation de $X$ dans $L$ soit finie
\'etale sur $X$.
\end{proposition}

\begin{proof}
Omis.
\end{proof}

\noindent
{\bf Changement de point base.} Pour tous points géométriques $\overline{x}_1, \overline{x}_2$
de $X$, il existe un isomorphisme des foncteurs fibres
$$
\mathcal{F}_{\overline{x}_1} \cong \mathcal{F}_{\overline{x}_2}
$$
(C'est un chemin de $\overline{x}_1$ à $\overline{x}_2$.) La conjugaison
par ce chemin donne un isomorphisme
$$
\pi_1(X, \overline{x}_1) \cong \pi_1(X, \overline{x}_2)
$$
bien défini à automorphisme intérieur près.

\medskip\noindent
{\bf Fonctorialité.} Pour tout morphisme $X_1\to X_2$ de schémas connexes
et tout $\overline{x}\in X_1$, il existe un morphisme canonique
$$
\pi_1(X_1, \overline{x}) \to \pi_1(X_2, \overline{x})
$$
(Pourquoi ? Parce que le foncteur fibre ...)

\medskip\noindent
{\bf Corps de base.} Soit $X$ une variété sur un corps $k$. On obtient alors
$$
\pi_1(X, \overline{x}) \to
\pi_1(Spec(k), \overline{x}) =^{\text{prop}} Gal(k^{\text{sep}}/k)
$$
Ce morphisme est surjectif si et seulement si $X$ est géométriquement connexe sur $k$.
Dans le cas géométriquement connexe, nous obtenons donc une suite exacte courte de groupes
profinis
$$
1 \to \pi_1(X_{\overline{k}}, \overline{x}) \to
\pi_1(X, \overline{x}) \to
Gal(k^{\text{sep}}/k) \to 1
$$
($\pi_1(X_{\overline{k}}, \overline{x})$ : groupe fondamental géométrique de
$X$, $\pi_1(X, \overline{x})$ : groupe fondamental arithmétique de $X$)

\medskip\noindent
{\bf Comparaison.} Si $X$ est une variété sur $\mathbf{C}$, alors
$$
\pi_1(X, \overline{x}) =
\text{ complété profini de }
\pi_1(X(\mathbf{C})(\text{ topologie usuelle}), x)
$$
(où $x\in X(\mathbf{C})$)

\medskip\noindent
{\bf Frobenius.} Soit $X$ une variété sur $k$, avec $\# k < \infty$. Pour tout $x \in X$
point fermé, soit
$$
F_x\in \pi_1(x, \overline{x}) =
\text{Gal}(\kappa(x)^{\text{sep}}/\kappa(x))
$$
le Frobenius géométrique.
Soit $\overline\eta$ un point géométrique algébrique au-dessus du point générique. Alors
$$
\pi_1(X, \overline\eta) \leftarrow^{\cong}
\pi_1(X, \overline{x})
{\text{fonctorialité} \atop \leftarrow} \pi_1(x, \overline{x})
$$

\noindent
Fait élémentaire :
$$
\begin{matrix}
\pi_1(X, \overline \eta) & \to^{\deg} \pi_1(\Spec(k), \overline \eta) * &
= Gal(k^{sep}/k) \\
& & || \\
& & \widehat{\mathbf{Z}}\cdot F_{\Spec(k)} \\
F_x & \mapsto & \deg(x)\cdot F_{\Spec(k)}
\end{matrix}
$$
Rappelons que $\deg(x) = [\kappa(x):k]$.

\medskip\noindent
{\bf Groupes fondamentaux et faisceaux lisses.}
Soit $X$ un schéma connexe et $\overline{x}$ un point géométrique. Il existe des
équivalences de catégories
\begin{center}
\begin{tabular}{@{}c@{}}
\((\Lambda\text{ anneau fini})\)\\[2pt]
\begin{tabular}{@{}c@{\quad$\leftrightarrow$\quad}c@{}}
\parbox[c]{0.38\textwidth}{\centering
faisceaux de \(\Lambda\)-modules localement constants constructibles sur
\(X_\etale\)}
&
\parbox[c]{0.38\textwidth}{\centering
modules finis (discrets) sur \(\Lambda\) munis d'une action continue de
\(\pi_1(X, \overline{x})\)}
\end{tabular}\\[10pt]
\((\ell\text{ premier})\)\\[2pt]
\begin{tabular}{@{}c@{\quad$\leftrightarrow$\quad}c@{}}
\parbox[c]{0.38\textwidth}{\centering
faisceaux \(\ell\)-adiques lisses}
&
\parbox[c]{0.38\textwidth}{\centering
un \(\mathbf{Z}_\ell\)-module \(M\) de type fini muni d'une action continue
de \(\pi_1(X, \overline{x})\) ; on utilise la topologie \(\ell\)-adique sur
\(M\)}
\end{tabular}
\end{tabular}
\end{center}
En particulier, les faisceaux lisses de $\mathbf{Q}_l$-modules correspondent aux
homomorphismes continus
$$
\pi_1(X, \overline{x}) \to \text{GL}_r(\mathbf{Q}_l), \quad r\geq 0
$$

\noindent
Notation : un module muni d'une action $(M, \rho)$ correspond au faisceau
$\mathcal{F}_\rho$.

\medskip\noindent
{\bf Formules des traces.} Soit $X$ une variété sur $k$, avec $\# k < \infty$.
\begin{enumerate}
\item $\Lambda$ est un anneau fini tel que $(\# \Lambda, \# k)=1$
$$
\rho : \pi_1(X, \overline{x})\to \text{GL}_r(\Lambda)
$$
continue. Pour tout $n\geq 1$, on a
$$
\sum_{d|n}d\left(
\sum_{x\in |X|, \atop \deg(x)=d}
\text{Tr}(\rho(F_x^{n/d}))\right) =
\text{Tr}\left(
(\pi_x^n)^* |_{R\Gamma_c(X_{\overline{k}}, \mathcal{F}_\rho)}\right)
$$
\item Soit $l\neq char(k)$ un nombre premier et soit $\rho : \pi_1(X, \overline{x})\to
\text{GL}_r(\mathbf{Q}_l)$. Pour tout $n\geq 1$,
$$
\sum_{d|n} d\left(
\sum_{x\in |X| \atop \deg(x)=d}
\text{Tr}
\left(
\rho(F_x^{n/d})
\right)
\right) =
\sum_{i = 0}^{2\dim X}
(-1)^i
\text{Tr}\left(
\pi_X^* |_{H_c^i(X_{\overline{k}}, \mathcal{F}_\rho)}\right)
$$
\end{enumerate}

\noindent
{\bf Conjectures de Weil.} (Deligne--Weil I, 1974) Si $X$ est lisse et projective sur $k$ et
$\# k = q$, alors les valeurs propres de $\pi_X^*$ sur $H^i(X_{\overline{k}},
\mathbf{Q}_l)$ sont des entiers algébriques $\alpha$ tels que $|\alpha|=q^{i/2}$.

\medskip\noindent
{\bf Conjectures de Deligne.} (presque entièrement démontrées par
Lafforgue et d'autres, $\ldots$) Soit $X$ une variété normale sur le corps fini $k$ et soit
$$
\rho : \pi_1(X, \overline{x}) \longrightarrow \text{GL}_r(\mathbf{Q}_l)
$$
continue. Supposons $\rho$ irréductible et $\det(\rho)$ d'ordre fini. Alors :
\begin{enumerate}
\item il existe un corps de nombres $E$ tel que, pour tout $x\in
|X|$ (point fermé), le polynôme caractéristique de $\rho(F_x)$ ait ses coefficients dans $E$ ;
\item pour tout $x\in |X|$, les valeurs propres $\alpha_{x, i}$, $i = 1, \ldots,
r$, de $\rho(F_x)$ ont toutes une valeur absolue complexe égale à $1$
(ce sont des nombres algébriques, pas nécessairement des entiers) ;
\item pour toute place finie $\lambda$ de $E$ ne divisant pas $p$
(quitte à agrandir légèrement $E$), il existe
$$
\rho\lambda : \pi_1(X, \overline{x}) \to \text{GL}_r(E_\lambda)
$$
compatible avec $\rho$ (les polynômes caractéristiques des $F_x$ sont les mêmes).
\end{enumerate}

\begin{theorem}[Deligne, Weil II]
\label{theorem-weil-II}
Pour un faisceau
$\mathcal{F}_\rho$ tel que $\rho$ satisfasse les conclusions de la conjecture
ci-dessus, les valeurs propres de $\pi_X^*$ sur $H_c^i(X_{\overline{k}},
\mathcal{F}_{\rho})$ sont des nombres algébriques $\alpha$ dont les valeurs absolues
$$
|\alpha|=q^{w/2}, \text{ pour }w\in \mathbf{Z},\ w\leq i
$$
De plus, si $X$ est lisse et projective, alors $w = i$.
\end{theorem}

\begin{proof}
Voir \cite{WeilII}.
\end{proof}




%12.03.09
\section{Groupes profinis, cohomologie et homologie}
\label{section-profinite-cohomology}

\noindent
Soit $G$ un groupe profini.

\medskip\noindent
{\bf Cohomologie.}
Considérons la catégorie des modules discrets munis d'une action continue de $G$.
Cette catégorie possède assez d'injectifs, et l'on peut définir
$$
H^i(G, M) = R^iH^0(G, M) = R^i(M\mapsto M^G)
$$
Il existe aussi une version dérivée $RH^0(G, -)$.

\medskip\noindent
{\bf Homologie.}
Considérons la catégorie des groupes abéliens compacts munis d'une action continue de $G$.
Cette catégorie possède assez de projectifs, et l'on peut définir
$$
H_i(G, M) = L_iH_0(G, M)=L_i(M\mapsto M_G)
$$
et il existe également une version dérivée.

\medskip\noindent
{\bf Dualité élémentaire.}
Le foncteur $M\mapsto M^\wedge = \Hom_{cont}(M, S^1)$
échange les catégories précédentes, et
$$
H^i(G, M)^\wedge = H_i(G, M^\wedge)
$$
De plus, ce foncteur envoie les $G$-modules discrets de torsion sur les
$G$-modules profinis continus, et réciproquement ; si $M$ est un module continu
discret ou profini muni d'une action de $G$, alors
$M^\wedge = \Hom(M, \mathbf{Q}/\mathbf{Z})$.

\medskip\noindent
{\bf Notes sur l'homologie.}
\begin{enumerate}
\item Si l'on considère les $\Lambda$-modules pour un anneau fini $\Lambda$,
alors on peut identifier
$$
H_i(G, M)=Tor_i^{\Lambda[[G]]}(M, \Lambda)
$$
où $\Lambda[[G]]$ est la limite des algèbres de groupe des quotients finis
de $G$.
\item Si $G$ est un sous-groupe normal de $\Gamma$, et si $\Gamma$ est également
profini, alors
\begin{itemize}
\item $H^0(G, -)$ : $\Gamma$-modules discrets$\to$ modules discrets
munis d'une action de $\Gamma/G$
\item $H_0(G, -)$ : $\Gamma$-modules compacts $\to$ modules compacts
munis d'une action de $\Gamma/G$
\end{itemize}
et le groupe profini $\Gamma/G$ agit donc sur les groupes de cohomologie
de $G$ à valeurs dans un $\Gamma$-module. Autrement dit, il existe des foncteurs
dérivés
$$
\begin{aligned}
RH^0(G, -) :\quad
&D^{+}(\text{modules discrets }\Gamma\text{-équivariants})\\
&\longrightarrow
D^{+}(\text{modules discrets }\Gamma/G\text{-équivariants})
\end{aligned}
$$
et de même pour $LH_0(G, -)$.
\end{enumerate}








\section{Cohomologie des courbes, revisitée}
\label{section-cohomology-curves-revisited}

\noindent
Soit $k$ un corps et soit $X$ une courbe lisse, géométriquement connexe sur $k$.
On dispose de la suite exacte fondamentale
$$
1 \to
\pi_1(X_{\overline{k}}, \overline \eta) \to
\pi_1(X, \overline\eta) \to
\text{Gal}(k^{^{sep}}/k) \to 1
$$
Si $\Lambda$ est un anneau fini tel que $\#\Lambda\in k^*$, si $M$ est un
$\Lambda$-module fini et si l'on se donne
$$
\rho : \pi_1(X, \overline\eta) \to \text{Aut}_{\Lambda}(M)
$$
continue, alors $\mathcal{F}_\rho$ désigne le faisceau associé sur
$X_\etale$.

\begin{lemma}
\label{lemma-identify-h2c}
Il existe un isomorphisme canonique
$$
H_c^2(X_{\overline{k}}, \mathcal{F}_\rho)=(M)_{\pi_1(X_{\overline{k}},
\overline\eta)}(-1)
$$
en tant que $\text{Gal}(k^{^{sep}}/k)$-modules.
\end{lemma}

\noindent
Ici, l'indice ${}_{\pi_1(X_{\overline{k}}, \overline\eta)}$
désigne les coinvariants, et $(-1)$ le twist à la Tate, c'est-à-dire que
$\sigma\in \text{Gal}(k^{^{sep}}/k)$ agit par
$$
\chi_{cycl}(\sigma)^{-1}.\sigma\text{ sur le membre de droite}
$$
où
$$
\chi_{cycl} :
\text{Gal}(k^{^{sep}}/k)
\to
\prod\nolimits_{l\neq char(k)}\mathbf{Z}_l^*
$$
est le caractère cyclotomique.

\medskip\noindent
Reformulation (Deligne, Weil II, page 338). Pour tout faisceau fini localement
constant $\mathcal{F}$ sur $X$, il existe un quotient maximal $\mathcal{F}\to
\mathcal{F}''$ tel que $\mathcal{F}''/X_{\overline{k}}$ soit un faisceau constant ; ainsi
$$
\mathcal{F}'' = (X\to \Spec(k))^{-1}F''
$$
où $F''$ est un faisceau sur $\Spec(k)$, c'est-à-dire un
$\text{Gal}(k^{^{sep}}/k)$-module. Alors
$$
H_c^2(X_{\overline{k}}, \mathcal{F})\to H_c^2(X_{\overline{k}},
\mathcal{F}'')\to F''(-1)
$$
est un isomorphisme.

\begin{proof}[Démonstration du lemme \ref{lemma-identify-h2c}]
Soit $Y\to^{\varphi}X$ le revêtement galoisien étale fini
correspondant à $\Ker(\rho) \subset \pi_1(X, \overline\eta)$. Ainsi
$$
\text{Aut}(Y/X)=Ind(\rho)
$$
est le groupe de Galois. Alors $\varphi^*\mathcal{F}_\rho =\underline M_Y$ et
$$
\varphi_*\varphi^*\mathcal{F}_\rho\to \mathcal{F}_\rho
$$
ce qui donne
\begin{align*}
& H_c^2(X_{\overline{k}}, \varphi_*\varphi^*\mathcal{F}_\rho) \to
H_c^2(X_{\overline{k}}, \mathcal{F}_\rho)\\
& =H_c^2(Y_{\overline{k}}, \varphi^*\mathcal{F}_\rho)\\
& =H_c^2(Y_{\overline{k}}, \underline M) = \oplus_{\text{composantes irréd.
de } \atop Y_{\overline{k}}}M
\end{align*}
$$
\Im(\rho) \to H_c^2(Y_{\overline{k}}, \underline M) =
\oplus_{\text{composantes irréd. de } \atop Y_{\overline{k}}}
M \to_{\Im(\rho) \text{équivariant}} H_c^2(X_{\overline{k}},
\mathcal{F}_{\rho}) \to^{\text{action }
\Im(\rho) \atop \text{triviale}}
$$
si $C/\overline{k}$ est une courbe irréductible, alors $H_c^2(C, \underline M)=M$.

\medskip\noindent
Puisque
$$
{\text{ensemble des composantes } \atop \text{irréductibles de }Y_k} =
\frac{Im(\rho)}{Im(\rho|_{\pi_1(X_{\overline{k}}, \overline \eta)})}
$$
On en conclut que $H_c^2(X_{\overline{k}}, \mathcal{F}_\rho)$ est un
quotient de $M_{\pi_1(X_{\overline{k}}, \overline \eta)}$. D'autre part,
il existe une surjection
$$
\mathcal{F}_\rho\to \mathcal{F}'' = {\text{ faisceau sur }
X\text{ associé à } \atop (M)_{\pi_1(X_{\overline{k}}, \overline
\eta)}\leftarrow\pi_1(X, \overline \eta)}
$$
$$
H_c^2(X_{\overline{k}}, \mathcal{F}_\rho)\to
M_{\pi_1(X_{\overline{k}}, \overline\eta)}
$$
L'action de Galois ainsi tordue provient du fait que
$H_c^2(X_{\overline{k}}, \mu_n)=^{\text{can}} \mathbf{Z}/n\mathbf{Z}$.
\end{proof}

\begin{remark}
\label{remark-projective}
On en conclut donc que, si $X$ est en outre projective,
on dispose, fonctoriellement en la représentation $\rho$,
des identifications
$$
H^0(X_{\overline{k}}, \mathcal{F}_\rho) =
M^{\pi_1(X_{\overline{k}}, \overline\eta)}
$$
et
$$
H_c^2(X_{\overline{k}}, \mathcal{F}_\rho) =
M_{\pi_1(X_{\overline{k}}, \overline \eta)}(-1)
$$
Bien entendu, si $X$ n'est pas projective, alors
$H^0_c(X_{\overline{k}}, \mathcal{F}_\rho) = 0$.
\end{remark}


\begin{proposition}
\label{proposition-curve-kpi1}
Soit $X/k$ comme précédemment, mais avec $X_{\overline{k}}\neq \mathbf{P}^1_{\overline{k}}$.
Les foncteurs
$
(M, \rho)\mapsto H_c^{2-i}(X_{\overline{k}}, \mathcal{F}_\rho)
$
sont les foncteurs dérivés à gauche de
$(M, \rho)\mapsto H_c^2(X_{\overline{k}}, \mathcal{F}_\rho)$
donc
$$
H_c^{2-i}(X_{\overline{k}}, \mathcal{F}_\rho) =
H_i(\pi_1(X_{\overline{k}}, \overline \eta), M)(-1)
$$
De plus, il existe une version dérivée, à savoir
$$
R\Gamma_c(X_{\overline{k}}, \mathcal{F}_\rho)
=
LH_0(\pi_1(X_{\overline{k}}, \overline \eta), M(-1))
=
M(-1)
\otimes_{\Lambda[[\pi_1(X_{\overline{k}}, \overline \eta)]]}^\mathbf{L}
\Lambda
$$
dans $D(\Lambda[[\widehat{\mathbf{Z}}]])$.
De même, les foncteurs
$(M, \rho)\mapsto H^i(X_{\overline{k}}, \mathcal{F}_\rho)$
sont les foncteurs dérivés à droite de
$(M, \rho)\mapsto M^{\pi_1(X_{\overline{k}}, \overline \eta)}$
donc
$$
H^i(X_{\overline{k}}, \mathcal{F}_\rho) =
H^i(\pi_1(X_{\overline{k}}, \overline \eta), M)
$$
De plus, il existe également une version dérivée dans ce cas.
\end{proposition}

\begin{proof}
(Idée) Montrer que les deux membres sont des $\delta$-foncteurs universels.
\end{proof}

\begin{remark}
\label{remark-poincare-groups}
La proposition et la dualité élémentaire donnent alors
$$
H^{2-i}_c(X_{\overline{k}}, \mathcal{F}_\rho)
\times
H^i(X_{\overline{k}}, \mathcal{F}_\rho^\wedge(1))
\to
\mathbf{Q}/\mathbf{Z}
$$
un accouplement parfait. Si $X$ est projective, c'est la dualité de Poincaré.
\end{remark}





\section{Formule des traces abstraite}
\label{section-abstract-trace-formula}

\noindent
Supposons donnée une extension de groupes profinis
$$
1 \to G \to \Gamma \xrightarrow{\deg} \widehat{\mathbf{Z}} \to 1
$$
On dit que $\Gamma$ {\it possède une formule des traces abstraite} si et seulement s'il
existe
\begin{enumerate}
\item un entier $q\geq 1$, et
\item pour tout $d\geq 1$, un ensemble fini $S_d$ et, pour chaque $x\in S_d$, une
classe de conjugaison $F_x \in \Gamma$ telle que $\deg(F_x) = d$
\end{enumerate}
de sorte que les conditions suivantes soient satisfaites :
\begin{enumerate}
\item pour tout $\ell$ ne divisant pas $q$, on a $\text{cd}_\ell(G)<\infty$, et
\item pour tout anneau fini $\Lambda$ tel que $q\in \Lambda^*$,
pour tout $\Lambda$-module projectif fini $M$ muni d'une action continue de
$\Gamma$, et pour tout $n > 0$, on a
$$
\sum\nolimits_{d|n}d \left(
\sum\nolimits_{x \in S_d}
\text{Tr}( F_x^{n/d} |_M)
\right)
=
q^n \text{Tr}(F^n|_{M \otimes_{\Lambda[[G]]}^{\mathbf{L}}\Lambda})
$$
dans $\Lambda^\natural$.
\end{enumerate}
Ici, $M \otimes_{\Lambda[[G]]}^{\mathbf{L}}\Lambda = LH_0(G, M)$ désigne
l'homologie dérivée, et $F=1$ dans $\Gamma/G = \widehat{\mathbf{Z}}$.

\begin{remark}
\label{remark-abstract-trace-formula}
Voici quelques observations au sujet de cette notion.
\begin{enumerate}
\item Pour modéliser les courbes projectives, on peut utiliser la cohomologie et
le facteur $q^n$ n'est pas nécessaire.
\item Les seuls exemples que je connaisse sont $\Gamma = \pi_1(X, \overline \eta)$,
où $X$ est lisse, géométriquement irréductible et un $K(\pi, 1)$ sur un corps
fini. Dans ce cas, $q = (\# k)^{\dim X}$. Compte tenu de la proposition, nous avons démontré
ce résultat pour les courbes dans ce cours.
\item Une fois l'entier $q$ donné, les ensembles $S_d$ sont déterminés de manière
unique. (On peut multiplier $q$ par un entier $m$, puis remplacer $S_d$ par
$m^d$ copies de $S_d$ sans changer la formule.)
\end{enumerate}
\end{remark}

\begin{example}
\label{example-commutative}
Fixons un entier $q\geq 1$.
$$
\begin{matrix}
1 &
\to &
G = \widehat{\mathbf{Z}}^{(q)} &
\to &
\Gamma &
\to &
\widehat{\mathbf{Z}} &
\to &
1 \\
&
&
= \prod_{l\not \mid q} \mathbf{Z}_l &
&
F &
\mapsto &
1
\end{matrix}
$$
avec $FxF^{-1} = ux$, où $u \in (\widehat{\mathbf{Z}}^{(q)})^*$.
En utilisant seulement les modules triviaux
$\mathbf{Z}/m\mathbf{Z}$, on voit que
$$
q^n - (qu)^n \equiv \sum\nolimits_{d|n} d\# S_d
$$
dans $\mathbf{Z}/m\mathbf{Z}$ pour tout $(m, q)=1$ (à
$u \to u^{-1}$ près) ; cela entraîne $qu = a\in \mathbf{Z}$
et $|a| < q$. Le cas particulier $a = 1$ se présente avec
$$
\Gamma = \pi_1^t(\mathbf{G}_{m, \mathbf{F}_p}, \overline \eta),
\quad
\# S_1 = q - 1,
\quad\text{et}\quad
\# S_2 = \frac{(q^2-1)-(q-1)}{2}
$$
\end{example}



%12.08.09
\section{Formes automorphes et faisceaux}
\label{section-automorphic}

\noindent
Références : voir notamment les remarquables articles
\cite{D1}, \cite{D2} et \cite{D0} de Drinfeld.

\medskip\noindent
{\bf Formes cuspidales non ramifiées.}
Soit $k$ un corps fini de caractéristique $p$.
Soit $X$ une courbe projective lisse géométriquement irréductible sur $k$.
Posons $K = k(X)$, le corps des fonctions de $X$.
Soit $v$ une place de $K$, ce qui revient à se donner un
point fermé $x\in X$. Soit $K_v$ le complété de $K$ en $v$ ;
c'est aussi le corps des fractions du complété de
l'anneau local de $X$ en $x$.
Notons $O_v\subset K_v$ l'anneau des entiers. Posons en outre
$$
O = \prod\nolimits_v O_v \subset \mathbf{A} = \prod_v' K_v
$$
et soit $\Lambda$ un anneau, avec $p$ inversible dans $\Lambda$.

\begin{definition}
\label{definition-unramified}
Une {\it forme cuspidale non ramifiée sur $\text{GL}_2(\mathbf{A})$ à valeurs dans
$\Lambda$}\footnote{Cette notation n'est probablement pas standard.}
est une fonction
$$
f : \text{GL}_2(\mathbf{A}) \to \Lambda
$$
telle que
\begin{enumerate}
\item $f(x\gamma) = f(x)$ pour tout $x\in \text{GL}_2(\mathbf{A})$ et tout
$\gamma\in \text{GL}_2(K)$
\item $f(ux) = f(x)$ pour tout $x\in \text{GL}_2(\mathbf{A})$ et tout
$u\in \text{GL}_2(O)$
\item pour tout $x\in \text{GL}_2(\mathbf{A})$,
$$
\int_{\mathbf{A} \mod K} f
\left(x
\left(
\begin{matrix}
1 & z \\
0 & 1
\end{matrix}
\right)
\right) dz = 0
$$
voir \cite[section 4.1]{dJ-conjecture}
pour une explication de la manière de donner un sens à cette expression
pour un anneau quelconque $\Lambda$ dans lequel $p$ est inversible.
\end{enumerate}
\end{definition}

\noindent
{\bf Opérateurs de Hecke.}
Pour une place $v$ de $K$ et une forme cuspidale non ramifiée $f$, posons
$$
T_v(f)(x) =
\int_{g\in M_v}f(g^{-1}x)dg,
$$
et
$$
U_v(f)(x) =
f\left(
\left(
\begin{matrix}
\pi_v^{-1} & 0 \\
0 & \pi_v^{-1}
\end{matrix}
\right)x\right)
$$
Notations : ici, $\pi_v \in O_v$ est une uniformisante,
$$
M_v =
\left\{
h\in Mat(2\times 2, O_v) | \det h = \pi_vO_v^*\right\}
$$
et $dg$ désigne la mesure de Haar sur $\text{GL}_2(K_v)$ telle que
$\int_{\text{GL}_2(O_v)} dg = 1$. Explicitement, on a
$$
T_v(f)(x) =
f\left(
\left(
\begin{matrix}
\pi_v^{-1}& 0 \\
0 & 1
\end{matrix}
\right)
x\right) +
\sum_{i = 1}^{q_v}
f\left(\left(
\begin{matrix}
1 & 0 \\
-\pi_v^{-1}\lambda_i
& \pi_v^{-1}
\end{matrix}
\right) x\right)
$$
où les $\lambda_i\in O_v$ forment un système de représentants de
$O_v/(\pi_v)=\kappa_v$, $q_v = \#\kappa_v$.

\medskip\noindent
{\bf Formes propres.} Une {\it forme propre} $f$ est une forme cuspidale non ramifiée
telle qu'une valeur de $f$ soit une unité, et telle que $T_vf = t_vf$ et
$U_vf = u_vf$ pour certains $t_v, u_v \in \Lambda$ déterminés de manière unique.

\begin{theorem}
\label{theorem-drinfeld-make-rho}
Étant donnée une forme propre $f$ à valeurs dans
$\overline{\mathbf{Q}}_l$, dont les valeurs propres
$u_v\in \overline{\mathbf{Z}}_l^*$, il existe une représentation
$$
\rho : \pi_1(X)\to \text{GL}_2(E)
$$
continue et absolument irréductible, où
$E$ est une extension finie de $\mathbf{Q}_\ell$ contenue dans
$\overline{\mathbf{Q}}_l$, et cette représentation vérifie
$t_v = \text{Tr}(\rho(F_v))$, et
$u_v = q_v^{-1}\det\left(\rho(F_v)\right)$ pour toute place $v$.
\end{theorem}

\begin{proof}
Voir \cite{D0}.
\end{proof}

\begin{theorem}
\label{theorem-drinfeld-make-f}
Supposons finie l'extension $\mathbf{Q}_l \subset E$, et soit
$$
\rho : \pi_1(X)\to \text{GL}_2(E)
$$
une représentation absolument irréductible et continue. Il existe alors une forme propre $f$ à
valeurs dans $\overline{\mathbf{Q}}_l$, dont les valeurs propres $t_v$, $u_v$
satisfont les égalités
$t_v = \text{Tr}(\rho(F_v))$ et $u_v = q_v^{-1}\det(\rho(F_v))$.
\end{theorem}

\begin{proof}
Voir \cite{D1}.
\end{proof}

\begin{remark}
\label{remark-lafforgue}
Grâce à Lafforgue et à de nombreux autres mathématiciens, nous disposons désormais
d'énoncés complets analogues aux deux théorèmes précédents pour $\text{GL}_n$,
y compris avec ramification !
Autrement dit, la correspondance de Langlands globale complète pour $\text{GL}_n$
est connue pour les corps de fonctions de courbes sur des corps finis. Cela
ne signifie toutefois pas qu'il ne reste pas de nombreuses questions intéressantes
sur les groupes fondamentaux des courbes sur les corps finis, comme
nous le verrons ci-dessous.
\end{remark}

\noindent
{\bf Caractère central.} Si $f$ est une forme propre, alors
$$
\begin{matrix}
\chi_f : &
O^*\backslash \mathbf{A}^*/K^* &
\to &
\Lambda^* \\
&
(1, \ldots, \pi_v, 1, \ldots, 1) &
\mapsto &
u_v^{-1}
\end{matrix}
$$
est appelé le {\it caractère central}. Il correspond, compte tenu des
normalisations ci-dessus, au déterminant de $\rho$. Posons
$$
C(\Lambda) =
\left\{
{\text{formes cuspidales non ramifiées } f \text{ à coefficients dans }\Lambda}
\atop {\text{ telles que } U_v f = \varphi_v^{-1}f\forall v}
\right\}
$$

\begin{proposition}
\label{proposition-cusp-forms-finite}
Si $\Lambda$ est noethérien, alors $C(\Lambda)$ est un
$\Lambda$-module de type fini. De plus, si $\Lambda$ est un corps et si
$\mathbf{F} \subset \Lambda$ est son sous-corps premier, alors
$$
C(\Lambda)=(C(\mathbf{F}))\otimes_{\mathbf{F}}\Lambda
$$
de manière compatible avec l'action de $T_v$.
\end{proposition}

\begin{proof}
Voir \cite[Proposition 4.7]{dJ-conjecture}.
\end{proof}

\noindent
Cette proposition entraîne immédiatement le lemme suivant.

\begin{lemma}
\label{lemma-eigenvalues-algebraic}
Algébricité des valeurs propres.
Si $\Lambda$ est un corps, alors les valeurs propres $t_v$ de $f\in
C(\Lambda)$ sont algébriques sur le sous-corps premier
$\mathbf{F} \subset \Lambda$.
\end{lemma}

\begin{proof}
Cela résulte de la proposition \ref{proposition-cusp-forms-finite}.
\end{proof}

\noindent
En combinant ce qui précède, on obtient l'astuce très utile suivante.

\begin{lemma}
\label{lemma-switch-l}
Changement de $l$. Soit $E$ un corps de nombres.
Partons d'une représentation
$$
\rho : \pi_1(X)\to SL_2(E_\lambda)
$$
absolument irréductible et continue, où $\lambda$ est une place de $E$
ne divisant pas $p$. Pour toute autre place $\lambda'$ de $E$
ne divisant pas $p$, il existe une extension finie $E'_{\lambda'}$
et une représentation continue absolument irréductible
$$
\rho': \pi_1(X)\to SL_2(E'_{\lambda'})
$$
compatible avec $\rho$, au sens où les polynômes caractéristiques
de tous les Frobenius sont les mêmes.
\end{lemma}

\noindent
Notons qu'il s'agit d'un cas particulier de la conjecture de Deligne !

\begin{proof}
Pour démontrer le lemme de changement, on applique
le théorème \ref{theorem-drinfeld-make-f}
pour obtenir une forme propre $f\in C(\overline{\mathbf{Q}}_l)$ associée à $\rho$.
On applique ensuite
la proposition \ref{proposition-cusp-forms-finite}
pour voir que l'on peut choisir $f\in C(E')$, où $E \subset E'$ est une extension finie.
En complétant ensuite $E'$, on obtient
une forme propre $f\in C(E'_{\lambda'})$, où
$E'_{\lambda'}$ est une extension finie de $E_{\lambda'}$.
Enfin, on applique
le théorème \ref{theorem-drinfeld-make-rho}
pour obtenir
$\rho': \pi_1(X) \to SL_2(E_{\lambda'}')$ absolument irréductible et continue,
quitte à agrandir encore légèrement $E'_{\lambda'}$.
\end{proof}

\noindent
Spéculation : si, pour un anneau (topologique) $\Lambda$, on a
$$
\left(
{\rho : \pi_1(X)\to SL_2(\Lambda) \atop \text{ absolument irréductible}}
\right)
\leftrightarrow
\text{ formes propres dans } C(\Lambda)
$$
alors toutes les valeurs propres de $\rho(F_v)$ sont algébriques (cet argument ne fonctionne
pas directement lorsque $\Lambda$ est un anneau fini). En partant de l'idée que la
correspondance de Langlands vaut plus généralement que pour les seuls corps,
on arrive à la conjecture suivante.

\medskip\noindent
{\bf Conjecture.}
(Voir \cite{dJ-conjecture}.)
Pour toute représentation continue
$$
\rho : \pi_1(X)\to \text{GL}_n(\mathbf{F}_l[[t]])
$$
on a $\# \rho(\pi_1(X_{\overline{k}}))<\infty$.

\medskip\noindent
Reformulation dans le langage des faisceaux :
« Pour tout faisceau lisse de $\overline{\mathbf{F}_l((t))}$-modules, la monodromie
géométrique est finie. »

\begin{theorem}
\label{theorem-conjecture-n-2}
La conjecture est vraie si $n\leq 2$.
\end{theorem}

\begin{proof}
Voir \cite{dJ-conjecture}.
\end{proof}

\begin{theorem}
\label{theorem-conjecture-l-bigger-2n}
La conjecture est vraie si $l > 2n$, sous réserve de certains résultats non démontrés.
\end{theorem}

\begin{proof}
Voir \cite{Gaitsgory}.
\end{proof}

\noindent
Il se trouve que la conjecture a une application.
Voir les travaux de Drinfeld sur les conjectures de Kashiwara. Il existe aussi
l'application beaucoup plus élémentaire suivante.

\begin{theorem}
\label{theorem-deformation-rings}
(Voir \cite[Théorème 3.5]{dJ-conjecture}.)
Supposons que
$$
\rho_0: \pi_1(X)\to \text{GL}_n(\mathbf{F}_l)
$$
soit continue, avec $l\neq p$. Supposons :
\begin{enumerate}
\item que la conjecture soit vraie pour $X$ ;
\item que $\rho_0 |_{\pi_1(X_{\overline{k}})}$ soit absolument irréductible ; et
\item que $l$ ne divise pas $n$.
\end{enumerate}
Alors l'anneau de déformation universel $R_{\text{univ}}$ de $\rho_0$ est
fini et plat sur $\mathbf{Z}_l$.
\end{theorem}

\noindent
Explication : il existe une représentation $\rho_{\text{univ}}:
\pi_1(X)\to \text{GL}_n(R_{\text{univ}})$ (anneau de déformation universel),
où $R_{\text{univ}}$ est local et
complet, de corps résiduel $\mathbf{F}_l$, et $(R_{\text{univ}}\to
\mathbf{F}_l)\circ\rho_{\text{univ}}\cong\rho_0$.
Étant donnés $R\to \mathbf{F}_l$, avec $R$ local et complet, et
$\rho : \pi_1(X)\to \text{GL}_n(R)$, il existe alors
$\psi : R_{\text{univ}}\to R$ tel que
$\psi\circ\rho_{\text{univ}}\cong \rho$. Le théorème affirme que le morphisme
$$
\Spec(R_{\text{univ}})
\longrightarrow
\Spec(\mathbf{Z}_l)
$$
est fini et plat. En particulier, une telle représentation $\rho_0$
se relève en une représentation $\rho : \pi_1(X) \to \text{GL}_n(\overline{\mathbf{Q}}_l)$.

\medskip\noindent
Notes :
\begin{enumerate}
\item Le théorème sur les déformations est facile.
\item Tout progrès vers la conjecture semble difficile.
\item Il serait intéressant de disposer de davantage de conjectures sur $\pi_1(X)$ !
\end{enumerate}




%12.10.09
\section{Comptage des points}
\label{section-counting}

\noindent
Soit $X$ une courbe lisse, géométriquement irréductible et
projective sur $k$, et posons $q = \# k$. La formule des traces donne ceci :
il existe des entiers algébriques $w_1, \ldots, w_{2g}$ tels que
$$
\# X(k_n) = q^n - \sum\nolimits_{i = 1}^{2g_X} w_i^n + 1.
$$
Si $\sigma\in \text{Aut}(X)$, alors, pour tout $i$, il existe $j$ tel que
$\sigma(w_i)=w_j$.

\medskip\noindent
{\bf Hypothèse de Riemann.} Pour tout $i$, on a $|\omega_i| = \sqrt{q}$.

\medskip\noindent
Cette hypothèse a été formulée par Emil Artin en 1924 pour les
courbes hyperelliptiques, puis démontrée par Weil en 1940. Weil en a donné deux démonstrations :
\begin{itemize}
\item l'une utilisant la théorie de l'intersection sur $X \times X$ et le
théorème de l'indice de Hodge ;
\item l'autre utilisant la jacobienne de $X$.
\end{itemize}
Il existe une autre démonstration, dont l'idée initiale est due à Stepanov et
que Bombieri a donnée : elle utilise le corps de fonctions $k(X)$ et
son opérateur de Frobenius (1969). Le point de départ est que, pour
$f\in k(X)$, la fonction rationnelle $f^q - f$
s'annule en tous les points $\mathbf{F}_q$-rationnels de $X$, et que l'on
peut essayer d'utiliser cette idée pour majorer le nombre de points.


\section{Forme précise du théorème de Chebotarev}
\label{section-chebotarev}

\noindent
Comme première application, démontrons une forme précise du théorème de Chebotarev
pour un revêtement galoisien étale fini de courbes.
Soit $\varphi : Y \to X$ un revêtement galoisien étale fini de
groupe $G$. Il correspond à un homomorphisme
$$
\pi_1(X) \longrightarrow G = \text{Aut}(Y/X)
$$
Supposons $Y_{\overline{k}}$ irréductible. Si $C\subset G$ est une classe de conjugaison,
alors, pour tout $n > 0$, on a
$$
| \# \{x \in X(k_n) \mid F_x \in C\} - \frac{\# C}{\# G} \cdot \# X(k_n) |
\leq
(\# C)(2g - 2) \sqrt{q^n}
$$
(Avertissement : vérifier soigneusement le coefficient $\# C$ au membre de droite
avant d'utiliser cette estimation.)

\begin{proof}[Esquisse]
Écrivons
$$
\varphi_*(\overline{\mathbf{Q}_l}) =
\oplus_{\pi \in \widehat{G}} \mathcal{F}_{\pi}
$$
où $\widehat{G}$ est l'ensemble des classes d'isomorphisme
des représentations irréductibles de
$G$ sur $\overline{\mathbf{Q}}_l$. Pour $\pi \in \widehat{G}$,
soit $\chi_{\pi}: G \to \overline{\mathbf{Q}}_l$
le caractère de $\pi$. Alors
$$
H^*(Y_{\overline{k}}, \overline{\mathbf{Q}}_l) =
\oplus_{\pi\in \widehat{G}}
H^*(Y_{\overline{k}}, \overline{\mathbf{Q}}_l)_\pi
=_{(\varphi\text{ finie })}
\oplus_{\pi\in \widehat{G}}
H^*(X_{\overline{k}}, \mathcal{F}_\pi)
$$
Si $\pi\neq 1$, alors
$$
H^0(X_{\overline{k}}, \mathcal{F}_\pi) =
H^2(X_{\overline{k}}, \mathcal{F}_\pi) = 0,\quad
\dim H^1(X_{\overline{k}}, \mathcal{F}_\pi) = (2g_X - 2)d_\pi^2
$$
(cela peut se déduire de la formule des traces en faisant agir sur ...) et l'on voit que
$$
|\sum_{x \in X(k_n)} \chi_\pi(\mathcal{F}_x)| \leq
(2g_X - 2) d_\pi^2\sqrt{q^n}
$$
Écrivons $1_C = \sum_\pi a_\pi \chi_\pi$ ; alors
$a_\pi = \langle 1_C, \chi_\pi\rangle$, et
$a_1 = \langle 1_C, \chi_1\rangle = \frac{\# C}{\# G}$, où
$$
\langle f, h\rangle = \frac{1}{\# G}\sum_{g \in G} f(g)\overline{h(g)}
$$
On a donc la relation
$$
\frac{\# C}{\# G} = ||1_C||^2 = \sum|a_\pi|^2
$$
Dernière étape :
\begin{align*}
\#\left\{x \in X(k_n) \mid F_x \in C\right\}
& =
\sum_{x \in X(k_n)} 1_C(x) \\
& =
\sum_{x \in X(k_n)} \sum_\pi a_\pi \chi_\pi(F_x) \\
& =
\underbrace{\frac{\# C}{\# G} \# X(k_n)}_{
\text{terme pour }\pi = 1}
+
\underbrace{\sum_{\pi\neq 1}a_\pi\sum_{x\in X(k_n)}\chi_\pi(F_x)}_{
\text{ terme d'erreur (à majorer par }E)}
\end{align*}
On peut majorer le terme d'erreur par
\begin{align*}
|E|
& \leq
\sum_{\pi \in \widehat{G}, \atop \pi \neq 1}
|a_\pi| (2g - 2) d_\pi^2 \sqrt{q^n} \\
& \leq
\sum_{\pi \neq 1} \frac{\# C}{\# G} (2g_X - 2) d_\pi^3 \sqrt{q^n}
\end{align*}
D'après la conjecture de Weil, $\# X(k_n)\sim q^n$.
\end{proof}



\section{Combien de nombres premiers se décomposent complètement ?}
\sectionmark{Nombres premiers complètement décomposés}
\label{section-how-many}

\noindent
Cette section donne une seconde application de l'hypothèse de Riemann pour les
courbes sur un corps fini. Les théoriciens des nombres pourront consulter
l'article d'Ihara intitulé
« Combien de nombres premiers se décomposent complètement dans une extension de Galois
non ramifiée infinie d'un corps global ? », voir \cite{Ihara}.
Considérons la suite exacte fondamentale
$$
1 \to
\pi_1(X_{\overline{k}}) \to
\pi_1(X) \xrightarrow{\deg}
\widehat{\mathbf{Z}} \to 1
$$

\begin{proposition}
\label{proposition-finite-set-frobenii-generate-topologically}
Il existe des points fermés $x_1, \ldots, x_n$ de $X$
tels que l'ensemble de {\bf tous} les éléments de Frobenius correspondant à ces
points engendre topologiquement $\pi_1(X)$.
\end{proposition}

\noindent
On peut aussi formuler l'énoncé ainsi :
il existe $x_1, \ldots, x_n\in |X|$ tels que
le plus petit sous-groupe normal fermé $\Gamma$ de $\pi_1(X)$
contenant exactement $1$ élément de Frobenius pour chaque $x_i$ soit égal à $\pi_1(X)$, c'est-à-dire
$\Gamma = \pi_1(X)$.

\begin{proof}
Choisissons $N\gg 0$ et posons
$$
\{x_1, \ldots, x_n\} =
{\text{ ensemble de tous les points fermés de}
\atop X \text{ de degré} \leq N\text{ sur } k}
$$
Soit $\Gamma\subset \pi_1(X)$ comme dans la reformulation précédente pour ces
points. Supposons $\Gamma \neq \pi_1(X)$. On peut alors choisir un sous-groupe ouvert normal
$U$ de $\pi_1(X)$ contenant $\Gamma$ et tel que
$U \neq \pi_1(X)$. Par l'hypothèse de Riemann pour $X$, notre ensemble de points contient un
$x_{i_1}$ de degré $N$ et un $x_{i_2}$ de degré $N - 1$. Cela montre que
$\deg : \Gamma \to \widehat{\mathbf{Z}}$ est surjective,
et il en va donc de même pour $U$. Cela signifie exactement que, si
$Y \to X$ est le revêtement galoisien étale fini
correspondant à $U$, alors $Y_{\overline{k}}$ est irréductible.
Posons $G = \text{Aut}(Y/X)$. On a
$$
Y \to^G X,\quad G = \pi_1(X)/U
$$
Par construction, tous les points de $X$ de degré $\leq N$ se décomposent
complètement dans $Y$. En particulier,
$$
\# Y(k_N)\geq (\# G)\# X(k_N)
$$
Appliquons l'hypothèse de Riemann aux deux membres. On obtient
$$
q^N+1+2g_Yq^{N/2}\geq \# G\# X(k_N)\geq \#
G(q^N+1-2g_Xq^{N/2})
$$
Comme $2g_Y-2 = (\# G)(2g_X-2)$, cela donne
$$
q^N + 1 + (\# G)(2g_X - 1) + 1)q^{N/2}\geq
\# G (q^N + 1 - 2g_Xq^{N/2})
$$
On voit donc que $G$ est nécessairement trivial si $N$ est assez grand.
\end{proof}

\noindent
{\bf Question curieuse.}
Posons $W_X = \deg^{-1}(\mathbf{Z})\subset \pi_1(X)$.
Existe-t-il un ensemble fini de points fermés $x_1, \ldots, x_n$ de $X$ tel que
l'ensemble de tous les éléments de Frobenius correspondant à ces points
engendre {\it algébriquement} $W_X$ ?

\medskip\noindent
Un argument de catégorie de Baire ramène cette question à la même question
pour tous les éléments de Frobenius.





\section{Combien y a-t-il réellement de points ?}
\label{section-really}

\noindent
Si le genre de la courbe est grand par rapport à $q$, alors le terme
principal de la formule $\# X(k) = q - \sum \omega_i + 1$ n'est pas le terme $q$,
mais le deuxième terme $\sum \omega_i$, qui peut avoir a priori une
taille de l'ordre de $2g_X\sqrt{q}$. Dans l'article \cite{Drinfeld-number},
Drinfeld et Vlăduţ montrent que ce maximum est en fait au plus de l'ordre de
$g\sqrt{q}$, comme Ihara l'avait prédit.

\medskip\noindent
Fixons $q$ et soit $k$ un corps à $q$ éléments. Posons
$$
A(q) = \limsup_{g_X \to \infty} \frac{\# X(k)}{g_X}
$$
où $X$ parcourt les courbes projectives lisses géométriquement irréductibles
sur $k$. Avec cette définition, on a les résultats suivants :
\begin{itemize}
\item HR $\Rightarrow A(q)\leq 2\sqrt{q}$
\item Ihara $\Rightarrow A(q)\leq \sqrt{2q}$
\item DV $\Rightarrow A(q)\leq \sqrt{q}-1$ (en fait, cette borne est atteinte si $q$
est un carré).
\end{itemize}

\begin{proof}
Étant donné $X$, soient $w_1, \ldots, w_{2g}$ comme précédemment et posons $g = g_X$.
Posons $\alpha_i = \frac{w_i}{\sqrt{q}}$ ; alors $|\alpha_i| = 1$. Si $\alpha_i$
apparaît, alors $\overline{\alpha}_i = \alpha_i^{-1}$ apparaît également. On a
$$
N = \# X(k) \leq X(k_r) = q^r + 1 - (\sum_i \alpha_i^r) q^{r/2}
$$
En réécrivant, on voit que, pour tout $r \geq 1$,
$$
-\sum_i \alpha_i^r \geq Nq^{-r/2} - q^{r/2} - q^{-r/2}
$$
Observons que
$$
0 \leq |\alpha_i^n +\alpha_i^{n-1} +\ldots +\alpha_i +1|^2
= (n + 1) + \sum_{j = 1}^n (n + 1 - j) (\alpha_i^j + \alpha_i^{-j})
$$
Ainsi,
\begin{align*}
2g(n+1) & \geq - \sum_i \left(\sum_{j = 1}^n (n+1-j)(\alpha_i^j
+\alpha_i^{-j})\right)\\
& =-\sum_{j = 1}^n (n+1-j)\left(\sum_i\alpha_i^j
+\sum_i\alpha_i^{-j}\right)
\end{align*}
En prenant la moitié, on obtient
\begin{align*}
g(n+1)& \geq - \sum_{j = 1}^n (n+1-j)(\sum_i\alpha_i^j)\\
& \geq N\sum_{j = 1}^n (n+1-j)q^{-j/2}-\sum_{j = 1}^n
(n+1-j)(q^{j/2}+q^{-j/2})
\end{align*}
On en déduit
$$
\frac{N}{g}\leq \left(\sum_{j = 1}^n \frac{n+1-j}{n+1}q^{-j/2} \right)^{-1}
\cdot
\left(
1 + \frac{1}{g} \sum_{j = 1}^n \frac{n + 1 - j}{n + 1}(q^{j/2} + q^{-j/2})
\right)
$$
Fixons $n$ et faisons tendre $g\to \infty$.
$$
A(q)\leq \left(\sum_{j = 1}^n \frac{n+1-j}{n+1}q^{-j/2}\right)^{-1}
$$
Donc
$$
A(q)\leq \lim_{n\to\infty}(\ldots) = \left(\sum_{j = 1}^\infty
q^{-j/2}\right)^{-1}=\sqrt{q}-1
$$
\end{proof}





\begin{multicols}{2}[\section{Autres chapitres}]
\noindent
Préliminaires
\begin{enumerate}
\item \hyperref[introduction-section-phantom]{Introduction}
\item \hyperref[conventions-section-phantom]{Conventions}
\item \hyperref[sets-section-phantom]{Théorie des ensembles}
\item \hyperref[categories-section-phantom]{Catégories}
\item \hyperref[topology-section-phantom]{Topologie}
\item \hyperref[sheaves-section-phantom]{Faisceaux sur les espaces}
\item \hyperref[sites-section-phantom]{Sites et faisceaux}
\item \hyperref[stacks-section-phantom]{Champs}
\item \hyperref[fields-section-phantom]{Corps}
\item \hyperref[algebra-section-phantom]{Algèbre commutative}
\item \hyperref[brauer-section-phantom]{Groupes de Brauer}
\item \hyperref[homology-section-phantom]{Algèbre homologique}
\item \hyperref[derived-section-phantom]{Catégories dérivées}
\item \hyperref[simplicial-section-phantom]{Méthodes simpliciales}
\item \hyperref[more-algebra-section-phantom]{Compléments d'algèbre}
\item \hyperref[smoothing-section-phantom]{Lissage des morphismes d'anneaux}
\item \hyperref[modules-section-phantom]{Faisceaux de modules}
\item \hyperref[sites-modules-section-phantom]{Modules sur les sites}
\item \hyperref[injectives-section-phantom]{Objets injectifs}
\item \hyperref[cohomology-section-phantom]{Cohomologie des faisceaux}
\item \hyperref[sites-cohomology-section-phantom]{Cohomologie sur les sites}
\item \hyperref[dga-section-phantom]{Algèbre différentielle graduée}
\item \hyperref[dpa-section-phantom]{Algèbre à puissances divisées}
\item \hyperref[sdga-section-phantom]{Faisceaux différentiels gradués}
\item \hyperref[hypercovering-section-phantom]{Hyperrecouvrements}
\end{enumerate}
Schémas
\begin{enumerate}
\setcounter{enumi}{25}
\item \hyperref[schemes-section-phantom]{Schémas}
\item \hyperref[constructions-section-phantom]{Constructions de schémas}
\item \hyperref[properties-section-phantom]{Propriétés des schémas}
\item \hyperref[morphisms-section-phantom]{Morphismes de schémas}
\item \hyperref[coherent-section-phantom]{Cohomologie des schémas}
\item \hyperref[divisors-section-phantom]{Diviseurs}
\item \hyperref[limits-section-phantom]{Limites de schémas}
\item \hyperref[varieties-section-phantom]{Variétés}
\item \hyperref[topologies-section-phantom]{Topologies sur les schémas}
\item \hyperref[descent-section-phantom]{Descente}
\item \hyperref[perfect-section-phantom]{Catégories dérivées de schémas}
\item \hyperref[more-morphisms-section-phantom]{Compléments sur les morphismes}
\item \hyperref[flat-section-phantom]{Compléments sur la platitude}
\item \hyperref[groupoids-section-phantom]{Schémas en groupoïdes}
\item \hyperref[more-groupoids-section-phantom]{Compléments sur les schémas en groupoïdes}
\item \hyperref[etale-section-phantom]{Morphismes étales de schémas}
\end{enumerate}
Thèmes de théorie des schémas
\begin{enumerate}
\setcounter{enumi}{41}
\item \hyperref[chow-section-phantom]{Homologie de Chow}
\item \hyperref[intersection-section-phantom]{Théorie de l'intersection}
\item \hyperref[pic-section-phantom]{Schémas de Picard des courbes}
\item \hyperref[weil-section-phantom]{Théories cohomologiques de Weil}
\item \hyperref[adequate-section-phantom]{Modules adéquats}
\item \hyperref[dualizing-section-phantom]{Complexes dualisants}
\item \hyperref[duality-section-phantom]{Dualité pour les schémas}
\item \hyperref[discriminant-section-phantom]{Discriminants et différentes}
\item \hyperref[derham-section-phantom]{Cohomologie de de Rham}
\item \hyperref[local-cohomology-section-phantom]{Cohomologie locale}
\item \hyperref[algebraization-section-phantom]{Géométrie algébrique et formelle}
\item \hyperref[curves-section-phantom]{Courbes algébriques}
\item \hyperref[resolve-section-phantom]{Résolution des surfaces}
\item \hyperref[models-section-phantom]{Réduction semi-stable}
\item \hyperref[functors-section-phantom]{Foncteurs et morphismes}
\item \hyperref[equiv-section-phantom]{Catégories dérivées de variétés}
\item \hyperref[pione-section-phantom]{Groupes fondamentaux de schémas}
\item \hyperref[etale-cohomology-section-phantom]{Cohomologie étale}
\item \hyperref[crystalline-section-phantom]{Cohomologie cristalline}
\item \hyperref[proetale-section-phantom]{Cohomologie pro-étale}
\item \hyperref[relative-cycles-section-phantom]{Cycles relatifs}
\item \hyperref[more-etale-section-phantom]{Compléments de cohomologie étale}
\item \hyperref[trace-section-phantom]{La formule des traces}
\end{enumerate}
Espaces algébriques
\begin{enumerate}
\setcounter{enumi}{64}
\item \hyperref[spaces-section-phantom]{Espaces algébriques}
\item \hyperref[spaces-properties-section-phantom]{Propriétés des espaces algébriques}
\item \hyperref[spaces-morphisms-section-phantom]{Morphismes d'espaces algébriques}
\item \hyperref[decent-spaces-section-phantom]{Espaces algébriques décents}
\item \hyperref[spaces-cohomology-section-phantom]{Cohomologie des espaces algébriques}
\item \hyperref[spaces-limits-section-phantom]{Limites d'espaces algébriques}
\item \hyperref[spaces-divisors-section-phantom]{Diviseurs sur les espaces algébriques}
\item \hyperref[spaces-over-fields-section-phantom]{Espaces algébriques sur des corps}
\item \hyperref[spaces-topologies-section-phantom]{Topologies sur les espaces algébriques}
\item \hyperref[spaces-descent-section-phantom]{Descente et espaces algébriques}
\item \hyperref[spaces-perfect-section-phantom]{Catégories dérivées d'espaces}
\item \hyperref[spaces-more-morphisms-section-phantom]{Compléments sur les morphismes d'espaces}
\item \hyperref[spaces-flat-section-phantom]{Platitude sur les espaces algébriques}
\item \hyperref[spaces-groupoids-section-phantom]{Groupoïdes dans les espaces algébriques}
\item \hyperref[spaces-more-groupoids-section-phantom]{Compléments sur les groupoïdes dans les espaces}
\item \hyperref[bootstrap-section-phantom]{Amorçage}
\item \hyperref[spaces-pushouts-section-phantom]{Sommes amalgamées d'espaces algébriques}
\end{enumerate}
Thèmes de géométrie
\begin{enumerate}
\setcounter{enumi}{81}
\item \hyperref[spaces-chow-section-phantom]{Groupes de Chow des espaces}
\item \hyperref[groupoids-quotients-section-phantom]{Quotients de groupoïdes}
\item \hyperref[spaces-more-cohomology-section-phantom]{Compléments sur la cohomologie des espaces}
\item \hyperref[spaces-simplicial-section-phantom]{Espaces simpliciaux}
\item \hyperref[spaces-duality-section-phantom]{Dualité pour les espaces}
\item \hyperref[formal-spaces-section-phantom]{Espaces algébriques formels}
\item \hyperref[restricted-section-phantom]{Algébrisation des espaces formels}
\item \hyperref[spaces-resolve-section-phantom]{Résolution des surfaces revisitée}
\end{enumerate}
Théorie des déformations
\begin{enumerate}
\setcounter{enumi}{89}
\item \hyperref[formal-defos-section-phantom]{Théorie formelle des déformations}
\item \hyperref[defos-section-phantom]{Théorie des déformations}
\item \hyperref[cotangent-section-phantom]{Le complexe cotangent}
\item \hyperref[examples-defos-section-phantom]{Problèmes de déformation}
\end{enumerate}
Champs algébriques
\begin{enumerate}
\setcounter{enumi}{93}
\item \hyperref[algebraic-section-phantom]{Champs algébriques}
\item \hyperref[examples-stacks-section-phantom]{Exemples de champs}
\item \hyperref[stacks-sheaves-section-phantom]{Faisceaux sur les champs algébriques}
\item \hyperref[criteria-section-phantom]{Critères de représentabilité}
\item \hyperref[artin-section-phantom]{Axiomes d'Artin}
\item \hyperref[quot-section-phantom]{Espaces de Quot et de Hilbert}
\item \hyperref[stacks-properties-section-phantom]{Propriétés des champs algébriques}
\item \hyperref[stacks-morphisms-section-phantom]{Morphismes de champs algébriques}
\item \hyperref[stacks-limits-section-phantom]{Limites de champs algébriques}
\item \hyperref[stacks-cohomology-section-phantom]{Cohomologie des champs algébriques}
\item \hyperref[stacks-perfect-section-phantom]{Catégories dérivées de champs}
\item \hyperref[stacks-introduction-section-phantom]{Introduction aux champs algébriques}
\item \hyperref[stacks-more-morphisms-section-phantom]{Compléments sur les morphismes de champs}
\item \hyperref[stacks-geometry-section-phantom]{La géométrie des champs}
\end{enumerate}
Thèmes de théorie des espaces de modules
\begin{enumerate}
\setcounter{enumi}{107}
\item \hyperref[moduli-section-phantom]{Champs de modules}
\item \hyperref[moduli-curves-section-phantom]{Espaces de modules de courbes}
\end{enumerate}
Divers
\begin{enumerate}
\setcounter{enumi}{109}
\item \hyperref[examples-section-phantom]{Exemples}
\item \hyperref[exercises-section-phantom]{Exercices}
\item \hyperref[guide-section-phantom]{Guide bibliographique}
\item \hyperref[desirables-section-phantom]{Desiderata}
\item \hyperref[coding-section-phantom]{Style de programmation}
\item \hyperref[obsolete-section-phantom]{Obsolète}
\item \hyperref[fdl-section-phantom]{Licence de documentation libre GNU}
\item \hyperref[index-section-phantom]{Index généré automatiquement}
\end{enumerate}
\end{multicols}


\bibliography{my}
\bibliographystyle{amsalpha}

\end{document}
